% concatenated from ParaFreeHP_Algs_FinalDraft.tex, my_preamble.tex
\documentclass{article}
\usepackage{graphicx} % Required for inserting images
\usepackage{amsthm}
\usepackage{amsmath}
\usepackage{amssymb}
\usepackage{graphicx}
\usepackage{color}
\usepackage{ifpdf}
\usepackage{url}
%\usepackage{hyperref}
\usepackage{algorithm}
%\usepackage{abstract}
\usepackage[usenames,dvipsnames]{xcolor}
%\usepackage{xfrac}
%\usepackage{comment}
\usepackage{paralist}
%\usepackage{natbib}

\usepackage{algorithm}
\usepackage{algpseudocode}
\usepackage{algorithmicx}
%\usepackage{algorithmic}
\usepackage[ruled,vlined, linesnumbered, algo2e]{algorithm2e}
\usepackage[ruled,vlined]{algorithm2e}


%\usepackage[sc]{mathpazo} % Use the Palatino font
\usepackage[T1]{fontenc} % Use 8-bit encoding that has 256 glyphs
\linespread{1.05} % Line spacing - Palatino needs more space between lines
%\usepackage{microtype} % Slightly tweak font spacing for aesthetics

%\usepackage[hmarginratio=1:1, top=1in, columnsep=20pt, margin=2.5cm]{geometry} % Document margins
\usepackage[letterpaper, margin=1.2in]{geometry}
\usepackage{multicol} % Used for the two-column layout of the document
\usepackage[hang, small,labelfont=bf,up,textfont=it,up]{caption} % Custom captions under/above floats in tables or figures
\usepackage{booktabs} % Horizontal rules in tables
\usepackage{float} % Required for tables and figures in the multi-column environment - they need to be placed in specific locations with the [H] (e.g. \begin{table}[H])
%\usepackage{hyperref} % For hyperlinks in the PDF

%\usepackage{lipsum} % Package to generate dummy text throughout this template

%\usepackage[sc]{mathpazo} % Use the Palatino font
%\usepackage[T1]{fontenc} % Use 8-bit encoding that has 256 glyphs
%\linespread{1.05} % Line spacing - Palatino needs more space between lines
%\usepackage{microtype} % Slightly tweak font spacing for aesthetics

%\usepackage[hmarginratio=1:1,top=32mm,columnsep=20pt]{geometry} % Document margins
%\usepackage[a4paper]{geometry}

%\usepackage{multicol} % Used for the two-column layout of the document
%\usepackage[hang, small,labelfont=bf,up,textfont=it,up]{caption} % Custom captions under/above floats in tables or figures
%\usepackage{booktabs} % Horizontal rules in tables
%\usepackage{float} % Required for tables and figures in the multi-column environment - they need to be placed in specific locations with the [H] (e.g. \begin{table}[H])
%\usepackage{hyperref} % For hyperlinks in the PDF

%\usepackage{lettrine} % The lettrine is the first enlarged letter at the beginning of the text
%\usepackage{paralist} % Used for the compactitem environment which makes bullet points with less space between them

%\usepackage{abstract} % Allows abstract customization
%\renewcommand{\abstractnamefont}{\normalfont\bfseries} % Set the "Abstract" text to bold
%\renewcommand{\abstracttextfont}{\normalfont\small}%\itshape} % Set the abstract itself to small italic text

%\usepackage{titlesec} % Allows customization of titles
%\renewcommand\thesection{\arabic{section}} % Roman numerals for the sections
%\renewcommand\thesubsection{\arabic{subsection}} % Roman numerals for subsections
%\titleformat{\section}[block]{\large\scshape\centering}{\thesection.}{0.5em}{} % Change the look of the section titles
%\titleformat{\subsection}[block]{\normalsize\bfseries}{\thesubsection.}{0.5em}{} % Change the look of the section titles

%+ Renewcommand
%\renewcommand{\thesection}{ \arabic{section}}
%\renewcommand{\thesubsection}{\arabic{section}.\arabic{subsection}}
%\renewcommand{\thesubsubsection}{\arabic{section}.\arabic{subsection}.\alph{subsubsection}}
\theoremstyle{plain}
\newtheorem{theorem}{Theorem}[section]
\newtheorem{corollary}{Corollary}[section]
\newtheorem{lemma}{Lemma}[section]
\newtheorem{proposition}{Proposition}[section]
\theoremstyle{definition}
\newtheorem{definition}{Definition}[section]
\newtheorem{assumption}{Assumption}[section]
%\theoremstyle{remark}
\newtheorem{remark}{Remark}[section]
% \newtheorem{algorithm}{\noindent\normalfont\textsc{Algorithm}}
%\usepackage[french,bordercolor=none,backgroundcolor=pink,linecolor=none,textsize=footnotesize,textwidth=0.75in,shadow]{todonotes}
\usepackage{todonotes}

%\usepackage{fancyhdr} % Headers and footers
%\pagestyle{fancy} % All pages have headers and footers
%\fancyhead{} % Blank out the default header
%\fancyfoot{} % Blank out the default footer
%%\fancyhead[C]{Technical Report} % Custom header text
%\fancyfoot[RO,LE]{\thepage} % Custom footer text
%
%\renewcommand{\headrulewidth}{0pt} %% to remove the line on top of the pages
%
%%+ Redefine the reference item.
%\makeatletter
%\renewcommand\@biblabel[1]{#1.}
%\makeatother

%%%%%%%%%%%%%%%%%%%%%%%%%%%%%%%%%%%%%%%%%%%%%
%+ End of definition of environment
%%%%%%%%%%%%%%%%%%%%%%%%%%%%%%%%%%%%%%%%%%%%%

%+ Renewcommand symbols
\newcommand{\ud}{\mathrm{d}} % differential sign
\newcommand{\Id}{\mathbb{I}}
\newcommand{\Pa}{\mathcal{P}}
\newcommand{\R}{\mathbb{R}}
\newcommand{\Hil}{\mathcal{H}}
\newcommand{\Rext}{\R\cup\{+\infty\}}
\newcommand{\barR}{[-\infty, \infty]}
\newcommand{\abs}[1]{\left\vert#1\right\vert}
\newcommand{\abss}[1]{\vert#1\vert}
\newcommand{\absc}[1]{\big\vert#1\big\vert}
\newcommand{\set}[1]{\left\{#1\right\}}
\newcommand{\sets}[1]{\{#1\}}
\newcommand{\norm}[1]{\left\Vert#1\right\Vert}
\newcommand{\norms}[1]{\Vert#1\Vert}
\newcommand{\tnorm}[1]{\vert\!\Vert#1\vert\!\Vert}
\newcommand{\tnorms}[1]{\vert\!\Vert#1\vert\!\Vert}
\newcommand{\epi}[1]{\mathrm{epi}\left(#1\right)}
\newcommand{\Eproof}{\hfill $\square$}
\newcommand{\prox}{\mathrm{prox}}
\newcommand{\tprox}[2]{\mathrm{prox}_{#1}\left(#2\right)}
\newcommand{\proj}{\mathrm{proj}}
\newcommand{\kproj}[2]{\mathrm{proj}_{#1}\left(#2\right)}
\newcommand{\relint}[1]{\mathrm{ri}\left(#1\right)}
\newcommand{\argmin}{\mathrm{arg}\!\displaystyle\min}
\newcommand{\argmax}{\mathrm{arg}\!\displaystyle\max}
\newcommand{\dom}[1]{\mathrm{dom}(#1)}
\newcommand{\Dom}[1]{\mathrm{Dom}(#1)}
\newcommand{\zero}[1]{{\boldsymbol{0}}}
\newcommand{\Exps}[2]{\mathbb{E}_{#1}\left[#2\right]}
\newcommand{\Expsk}[2]{\mathbb{E}_{#1}\big[#2\big]}
\newcommand{\Expk}[1]{\mathbb{E}\big[#1\big]}
\newcommand{\Expsb}[2]{\mathbb{E}_{#1}\big[#2\big]}
\newcommand{\Prob}[1]{\mathbf{Prob}\left(#1\right)}
\newcommand{\cProb}[2]{\mathbf{Prob}\left(#1 \mid #2 \right)}
\newcommand{\zer}[1]{\mathrm{zer}(#1)}
\newcommand{\graph}[1]{\mathrm{graph}(#1)}
\newcommand{\gra}[1]{\mathrm{gra}(#1)}
\newcommand{\range}[1]{\mathrm{range}(#1)}
\newcommand{\mcal}[1]{\mathcal{#1}}
%\newcommand{\mbf}[1]{\mathbf{#1}}

\newcommand{\xopt}{x^{\star}}
\newcommand{\xbar}{\bar{x}}
\newcommand{\xast}{x^{\ast}}
\newcommand{\yast}{y^{\ast}}
\newcommand{\xb}{x}
\newcommand{\yb}{y}
\newcommand{\eb}{\mathbf{e}}
\newcommand{\zb}{z}
\newcommand{\wb}{w}
\newcommand{\ub}{u}
\newcommand{\vb}{v}
\newcommand{\lb}{\mathbf{l}}
\renewcommand{\sb}{s}
\newcommand{\rb}{r}
\newcommand{\qb}{q}
\newcommand{\ab}{a}
\newcommand{\bb}{b}
\newcommand{\cb}{c}
\newcommand{\db}{d}
\newcommand{\hb}{h}
\newcommand{\nb}{n}
\newcommand{\pb}{p}
\newcommand{\Ab}{\mathbf{A}}
\newcommand{\Gb}{\mathbf{G}}
\newcommand{\Hb}{\mathbf{H}}
\newcommand{\Kb}{\mathbf{K}}
\newcommand{\Ac}{\mathcal{A}}
\newcommand{\Hc}{\mathcal{H}}
\newcommand{\Bb}{\mathbf{B}}
\newcommand{\Db}{\mathbf{D}}
\newcommand{\Pb}{\mathbf{P}}
\newcommand{\Pbb}{\mathbb{P}}
\newcommand{\Qb}{\mathbf{Q}}
\newcommand{\Bc}{\mathcal{B}}
\newcommand{\Cb}{\mathbf{C}}
\newcommand{\Ub}{\mathbf{U}}
\newcommand{\Vb}{\mathbf{V}}
\newcommand{\Xb}{\mathbf{X}}
\newcommand{\Yb}{\mathbf{Y}}
\newcommand{\Zb}{\mathbf{Z}}
\newcommand{\Xc}{\mathcal{X}}
\newcommand{\Yc}{\mathcal{Y}}
\newcommand{\Zc}{\mathcal{Z}}
\newcommand{\Sc}{\mathcal{S}}
\newcommand{\Sdp}{\mathbb{S}}
\newcommand{\Dc}{\mathcal{D}}
\newcommand{\Lc}{\mathcal{L}}
\newcommand{\Qc}{\mathcal{Q}}
\newcommand{\Gc}{\mathcal{G}}
\newcommand{\Mc}{\mathcal{M}}
\newcommand{\Tc}{\mathcal{T}}
\newcommand{\Lb}{\mathbf{L}}
\newcommand{\Fb}{\mathbf{F}}
\newcommand{\Ib}{\mathcal{I}}
\newcommand{\Fc}{\mathcal{F}}
\newcommand{\Uc}{\mathcal{U}}
\newcommand{\Vc}{\mathcal{V}}
\newcommand{\Cc}{\mathcal{C}}
\newcommand{\Ec}{\mathcal{E}}
\newcommand{\Nc}{\mathcal{N}}
\newcommand{\Pc}{\mathcal{P}}
\newcommand{\Kc}{\mathcal{K}}
\newcommand{\Wc}{\mathcal{W}}
\newcommand{\Xopt}{\mathcal{X}^{\star}}
\newcommand{\Yopt}{\mathcal{Y}^{\star}}
\newcommand{\Lag}{\widetilde{\mathcal{L}}}
\newcommand{\iprod}[1]{\left\langle #1\right\rangle}
\newcommand{\iprods}[1]{\langle #1\rangle}
%\newcommand{\todo}[1]{\textbf{\color{red}Todo:}{\color{blue}~#1}}
\newcommand{\Exp}[1]{\mathbb{E}\left[#1\right]}
\newcommand{\cExp}[2]{\mathbb{E}\left[#1\mid #2\right]}
\newcommand{\Var}[1]{\mathrm{Var}\left\{#1\right\}}
\newcommand{\intx}[1]{\mathrm{int}\left(#1\right)}
\newcommand{\ri}[1]{\mathrm{ri}\left(#1\right)}
\renewcommand{\vec}[1]{\mathrm{vec}(#1)}
\newcommand{\mat}[1]{\mathrm{mat}(#1)}
\newcommand{\cond}[1]{\mathrm{cond}(#1)}
\newcommand{\rank}[1]{\mathrm{rank}(#1)}
%\newcommand{\diag}[1]{\mathrm{diag}(#1)}
\newcommand{\0}{\boldsymbol{0}}
\newcommand{\trace}[1]{\mathrm{trace}\left(#1\right)}
\newcommand{\kron}{\otimes}
\newcommand{\dist}[1]{\mathrm{dist}\left(#1\right)}
\newcommand{\kdist}[2]{\mathrm{dist}_{#1}\left(#2\right)}
\newcommand{\diag}[1]{\mathrm{diag}\left(#1\right)}
\newcommand{\gsc}{\textit{generalized self-concordant}~}
\newcommand{\Fgsc}{\mathcal{F}^{\mathrm{gsc}}_{M_f,\alpha}}
\newcommand{\SFunc}[2]{\mathcal{C}^{#1}\left(#2\right)}
\newcommand{\BigO}[1]{\mathcal{O}\left(#1\right)}
\newcommand{\BigOs}[1]{\mathcal{O}\big(#1\big)}
\newcommand{\SmallOs}[1]{o\big(#1\big)}
\newcommand{\SmallO}[1]{o\left(#1\right)}
\newcommand{\Ball}[2]{\mathbb{B}_{#1}\left(#2\right)}
\newcommand{\BigOc}[1]{{\color{blue}\mathcal{O}\left(#1\right)}}
\newcommand{\SmallOc}[1]{{\color{blue}o\left(#1\right)}}
\newcommand{\Up}[1]{\mathbb{U}_{\mathbf{\hat{q}}}\left(#1\right)}
\newcommand{\Uq}[1]{\mathbb{U}_{\mathbf{q}}\left(#1\right)}

%% For text colors
\newcommand{\mred}[1]{\textcolor{red}{#1}}
\newcommand{\mblue}[1]{\textcolor{blue}{#1}}

% Editing commands
\newcommand{\question}[1]{\noindent\textbf{\color{blue}Question: }{\raggedright\textcolor{BurntOrange}{#1}}}
\newcommand{\quocmargin}[1]{\marginpar{\raggedright\textcolor{BurntOrange}{#1}}}
%\newcommand{\quoctd}[1]{\textbf{\color{red}Quoc: }\textcolor{BurntOrange}{#1}}
\newcommand{\quoctd}[1]{\textbf{\color{red}Quoc: }\textcolor{blue}{#1}}
\newcommand{\nhanp}[1]{\textbf{\color{purple}Nhan: }\textcolor[rgb]{0.1,0.1,0.8}{#1}}
\newcommand{\lamn}[1]{\textbf{\color{magenta}Lam: }\textcolor[rgb]{0.8,0.1,0.1}{#1}}
\newcommand{\corr}[1]{{\color{green}#1}}
\newcommand{\mnote}[1]{{\color{magenta}#1}}
\newcommand{\mytbi}[1]{\textbf{\textit{#1}}}
\newcommand{\mytb}[1]{\textbf{#1}}
\newcommand{\myti}[1]{\textit{#1}}
\newcommand{\mbf}[1]{\mathbf{#1}}
\newcommand{\mbb}[1]{\mathbb{#1}}

% Gap between section/subsection.
\newcommand{\beforesubsec}{\vspace{-1.5ex}}
\newcommand{\aftersubsec}{\vspace{-1ex}}
\newcommand{\beforesec}{\vspace{-1.25ex}}
\newcommand{\aftersec}{\vspace{-1ex}}
\newcommand{\beforesubsubsec}{\vspace{-1.5ex}}
\newcommand{\aftersubsubsec}{\vspace{-1.5ex}}
\newcommand{\beforepara}{\vspace{-2.5ex}}


\newcommand{\myeq}[2]{%\vspace{-0.5ex}
\begin{equation}\label{#1}
{#2}
%\vspace{-0.75ex}
\end{equation}
}
\newcommand{\myeqn}[1]{%\vspace{-0.5ex}
\begin{equation*}
{#1}
%\vspace{-0.75ex}
\end{equation*}
}



\usepackage{algorithm,hyperref}
\usepackage{algpseudocode}
\usepackage{float}    % for [!htbp]
\usepackage{booktabs} % For professional quality tables (\toprule, \midrule, \bottomrule)
\usepackage{siunitx}  % For aligned numerical columns and scientific notation
\usepackage{wrapfig}

\title{Provable Parameter-Free Fixed-Point Algorithms with Linear Convergence Rates}
\author{
Quoc Tran-Dinh$^{*}$, Pham Ngoc Anh$^{\dagger}$, and Ha Manh Tien$^{\ddagger}$\\
$^{*}$Department of Statistics and Operations Research\\The University of North Carolina at Chapel Hill, USA; \\
$^{\dagger}$Laboratory of Applied Mathematics and Computing, \\
Posts and Telecommunications Institute of Technology, Hanoi, Vietnam;\\
$^{\ddagger}$Faculty of Basic Sciences and Foreign Languages,\\
Fire and Rescue Academy, Hanoi, Vietnam.
}
\date{June 2026}

\begin{document}

\maketitle
\begin{abstract}
In this paper, we develop provable parameter-free and adaptive fixed-point algorithms for contractive mappings, with an emphasis on automatically exploiting hidden contractivity without requiring prior knowledge of the contraction factor. 
Our first method is a completely parameter-free variant of the Halpern fixed-point iteration. 
It requires no line search, bisection, or prior estimate of the contraction factor, while retaining essentially the same per-iteration computational cost as classical fixed-point schemes. 
We establish explicit linear convergence rates for both the fixed-point residual and the distance to the unique fixed point. 
The second algorithm is an adaptive Halpern method that requires only an upper bound on the contraction factor and reduces to an existing adaptive Halpern scheme in the nonexpansive case. 
This method also enjoys explicit linear convergence guarantees. 
We further extend these ideas in two directions. 
First, by combining the proposed fixed-point schemes with Tikhonov regularization, we obtain a parameter-free method for solving co-coercive equations and establish an iteration complexity of $\mathcal{O}({\epsilon^{-1}\ln(\epsilon^{-1})})$ for computing an $\epsilon$-solution. 
Second, using the relation between Halpern iterations and Nesterov's accelerated fixed-point schemes, we derive parameter-free Nesterov's accelerated variants that inherit linear convergence in the contractive setting.
Numerical experiments on several examples demonstrate that the proposed algorithms are competitive with, and often outperform, existing adaptive fixed-point methods. In particular, the methods successfully exploit contractive behavior when it is present while remaining effective on nonexpansive problems. 
\end{abstract}

%%% 1. Introduction.
\section{Introduction}\label{sec:intro}
Developing parameter-free algorithms for optimization and related problems such as minimax optimization and variational inequalities has attracted renewed interest in recent years, largely due to their strong empirical performance across a broad range of applications, particularly in machine learning and data science \cite{Barzilai1988,carmon2022making,defazio2022parameter,defazio2023learning,Duchi2011,fercoq2019adaptive,ito2023parameter,KingmaB14,lan2026optimal,li2025simple,malitsky2020adaptive,orabona2016coin}. 
In contrast, the development of parameter-free methods for fixed-point problems has received considerably less attention. 
Although several parameter-free fixed-point algorithms have recently been proposed, their practical performance has been investigated only to a limited extent, with relatively few numerical studies reported in the literature (see, e.g., \cite{alakoya2021modified,diakonikolas2020halpern,lv2026preconditioned,malitsky2020forward,ogwo2025inertial,shen2026parameter,tan2022self}). 
Moreover, many of these methods still lack rigorous theoretical guarantees. 
Most existing results establish only asymptotic convergence, without providing explicit convergence rates or iteration-complexity estimates. 
This noticeable gap between the optimization and fixed-point literature motivates the present work.

%%% (a) Our motivation and goals.
\vspace{0.75ex}
\noindent\textbf{$\mathrm{(a)}$~Our motivation and research question.}
Apart from the growing success of parameter-free methods in optimization, our second motivation stems from a practical observation. 
In many applications of fixed-point methods, the underlying mapping is simply assumed to be nonexpansive, and algorithms designed for nonexpansive mappings are subsequently employed to approximate its fixed point. 
In practice, however, it is often difficult, and sometimes impossible, to determine whether the mapping is actually contractive, even locally. 
If the mapping is indeed contractive, then one would naturally expect an algorithm to enjoy a linear convergence rate. 
In contrast, for general nonexpansive mappings, the best possible convergence rate is typically no better than $\BigOs{1/k}$, which is significantly slower than linear convergence. 
This naturally leads to the following question:
\begin{center}
\textit{\textbf{Can we develop efficient parameter-free algorithms that automatically exploit hidden contractivity without requiring prior knowledge of the contraction factor?}}
\end{center}
%%%
To further investigate this question, we conduct a preliminary experiment on two representative mappings: (i) a linear mapping $T(x)=Q_1x+q_1$ and (ii) a nonlinear mapping $T(x)=c\arctan(Q_2x+q_2)$, where $c=0.99$, and $(Q_1,q_1)$ and $(Q_2,q_2)$ are chosen such that $T$ is $\rho$-contractive for some $\rho\in(0,1)$, but $\rho$ is close to one. 
We compare the seven adaptive algorithms proposed in this paper with Algorithm~3.1 of \cite{he2024convergence}, which is designed for nonexpansive mappings, and the geometric Halpern fixed-point method of \cite{park2022exact}, which assumes that the exact contraction factor $\rho$ is known.

Figure~\ref{fig:paired_convergence} plots the relative residual norms against the number of iterations. 
Here, the geometric Halpern method refers to the algorithm in \cite{park2022exact}, which uses the exact value of the contraction factor $\rho$.

\begin{figure}[!h]
	\centering
	\includegraphics[width=\textwidth]{figs/example_5_convergence}
	\caption{Convergence behavior of the nine algorithms for approximating a fixed point of $T$ on linear and nonlinear contractive mappings.}
	\label{fig:paired_convergence}
\end{figure}

As illustrated in Figure~\ref{fig:paired_convergence}, all seven of our proposed methods exhibit a clear linear convergence rate on both the linear and nonlinear test problems, reaching a relative residual tolerance of $10^{-8}$. 
The \texttt{Geo.~HP} is a Halpern fixed-point scheme using known contractive parameter $\rho$ also showed its linear convergence rate as theoretically stated in \cite{park2022exact}.
In contrast, Algorithm~3.1 of \cite{he2024convergence} initially displays a similar linear convergence trend but slows down considerably once the relative residual reaches approximately $10^{-4}$. 
These experiments indicate that our methods successfully exploit the hidden contractivity of the underlying mapping, thereby maintaining the expected linear convergence behavior while continuing to make steady progress toward high-accuracy solutions.

One may argue that, whenever $T$ is contractive, it is sufficient to apply the classical Banach-Picard (BP) iteration, which is guaranteed to converge linearly to the unique fixed point. 
While this is true under the assumption of global contractivity, such an assumption is often difficult to verify in practice. 
In many applications, it is unclear whether the underlying mapping is globally contractive, only locally contractive, or merely nonexpansive. 
Consequently, the BP iteration may fail to converge when initialized outside the region where contractivity holds. 
By contrast, the Halpern and Krasnosel'ski\v{\i}-Mann (KM) fixed-point iterations remain convergent under the much weaker assumption that $T$ is nonexpansive, making them substantially more robust in situations where the contractive behavior of $T$ is unknown. 
This observation motivates our focus on developing parameter-free Halpern- and KM-type methods that automatically exploit hidden contractivity whenever it is present while retaining the robustness of algorithms designed for nonexpansive mappings.

\vspace{0.75ex}
\noindent\textbf{$\mathrm{(b)}$~Problem statements.}
Let $T : \Hil \to \Hil$ be a mapping defined on a Hilbert space $\Hil$.
We are interested in approximating a fixed point of $T$, namely:
\begin{equation}\label{eq:FP}
\textrm{Find $x^{\star}\in \Hil$ such that:}~ x^{\star} = T(x^{\star}).
\tag{FP}
\end{equation}
Denote by $\mathrm{Fix}(T) := \sets{x^{\star} \in \Hil : x^{\star} = T(x^{\star}) }$ the set of fixed points of $T$.

Alternatively, we also consider the following root-finding problem:
\begin{equation}\label{eq:CE}
\textrm{Find $x^{\star}\in \Hil$ such that:}~ G(x^{\star}) = 0.
\tag{CE}
\end{equation}
where $G : \Hil \to \Hil$ is a single-valued mapping.
Denote by $\mathrm{zer}(G) := \sets{x^{\star} \in \Hil : G(x^{\star}) = 0}$ the solution set of \eqref{eq:CE}.
Throughout this paper, we assume that both $\mathrm{Fix}(T)$ and $\mathrm{zer}(G)$ are nonempty.

Problems \eqref{eq:FP} and \eqref{eq:CE} are closely related. Indeed, if $x^{\star} \in \mathrm{Fix}(T)$, then $x^{\star} \in \mathrm{zer}(\Id - T)$. 
Conversely, if $x^{\star} \in \mathrm{zer}(G)$, then $x^{\star} \in \mathrm{Fix}(\Id - \eta G)$ for any $\eta > 0$. 
Consequently, the fixed-point problem \eqref{eq:FP} can be viewed as an equivalent form of the root-finding problem \eqref{eq:CE}.

To develop numerical methods for solving \eqref{eq:FP}, it is common to assume that $T$ is either nonexpansive or $\rho$-contractive for some $\rho \in (0,1)$, see Section~\ref{sec:prelim_results} for definitions. 
Throughout this paper, we unify these two settings by referring to $T$ as a $\rho$-contractive mapping with $\rho \in (0,1]$, where the case $\rho=1$ corresponds to nonexpansive mappings. 
For the root-finding problem \eqref{eq:CE}, a standard assumption is that $G$ is $\beta$-co-coercive, i.e., $\iprods{G(x) - G(y), x - y} \geq \beta\norms{G(x) - G(y)}^2$ for all $x, y\in\Hil$.
In this case, we refer to \eqref{eq:CE} as a co-coercive equation.

The notions of nonexpansiveness and cocoercivity are closely connected. 
It is well-known that if $G$ is $\beta$-co-coercive, then $T := \Id - 2\beta G$ is nonexpansive. 
Conversely, if $T$ is nonexpansive, then $G := \Id - T$ is $\frac{1}{2}$-co-coercive. 
Therefore, solving a co-coercive equation is equivalent to approximating a fixed point of a nonexpansive mapping (or more generally, an averaged mapping), see \cite{Bauschke2011}.

For the fixed-point problem~\eqref{eq:FP}, if $T$ is globally contractive, i.e., $\norms{T(x) - T(y)} \leq \rho\norms{x - y}$ for all $x, y \in \Hil$ with $\rho \in (0, 1)$, then the classical Banach--Picard (BP) iteration \cite{Banach1922} converges linearly to the unique fixed point. 
When $T$ is only nonexpansive, two classical fixed-point methods are commonly employed. 
The first is the Krasnosel'ski\v{\i}-Mann (KM) iteration \cite{Krasnoselskii1955,mann1953mean}, which achieves the sublinear convergence rate $\BigOs{1/\sqrt{k}}$ in terms of the residual norm $\norms{x^k - T(x^k)}$. 
The second is the Halpern fixed-point iteration \cite{halpern1967fixed}, which is the primary focus of this paper.

For convenience, we first briefly recall the Halpern fixed-point iteration from \cite{halpern1967fixed}.
Starting from an initial point $x^0 \in \Hil$ and an anchor point $u\in \Hil$ (which may coincide with $x^0$), at each iteration $k \geq 0$, the Halpern fixed-point iteration updates
\begin{equation}\label{eq:HP_iteration_rule}
x^{k+1} = \lambda_k u + (1-\lambda_k)T(x^k),
\end{equation}
where $\lambda_k \in (0, 1)$ is a relaxation parameter. Various choices of $\lambda_k$ have been proposed to guarantee the convergence of \eqref{eq:HP_iteration_rule}, see, e.g., \cite{he2024convergence}. 
More recently, \cite{lieder2021convergence,sabach2017first} showed that the choice $\lambda_k = \frac{1}{k+2}$ yields the non-asymptotic convergence rate $\BigOs{1/k}$ for the residual norm $\norms{x^k - T(x^k)}$. 
Furthermore, this rate is known to be optimal, see, for example, \cite{lieder2021convergence,park2022exact}.

The development of parameter-free and adaptive Halpern methods is relatively recent. 
Diakonikolas \cite{diakonikolas2020halpern} appears to be among the first to propose a parameter-free variant of \eqref{eq:HP_iteration_rule} for solving the co-coercive equation \eqref{eq:CE}. 
However, the proposed method still relies on a line-search procedure usually used in optimization and VIPs. 
More recently, He et al.~\cite{he2024convergence} proposed the following simple adaptive update rule for $\lambda_k$ in \eqref{eq:HP_iteration_rule}, which avoids both line search and bisection:
\begin{equation}\label{eq:adaptive_stepsize}
\lambda_k = \frac{1}{1+\varphi_k}, \quad\text{where} \quad \varphi_k = 1 + \frac{2\iprods{x^{k-1} - T(x^{k-1}), x^0 - x^{k-1}}}{\norms{x^{k-1} - T(x^{k-1})}^2}.
\end{equation}
Under this update rule, the authors established a non-asymptotic convergence rate of $\BigOs{1/k}$ for the residual norm and proved that the sequence $\sets{x^k}$ converges strongly to the projection of $x^0$ onto the fixed-point set. 
However, this adaptive rule is intentionally and specifically designed for nonexpansive mappings and therefore still achieves only a sublinear convergence rate.

Motivated by and complementing the work of \cite{diakonikolas2020halpern,he2024convergence}, we develop in this paper two new algorithms for approximating a fixed point of \eqref{eq:FP} when $T$ is assumed to be $\rho$-contractive, though we do not need to know $\rho$. 
The first algorithm is inspired by the geometric Halpern method of \cite{park2022exact}, while the second builds upon the adaptive parameter update introduced in \cite{he2024convergence}. 
Both methods are designed to exploit the contractivity of $T$ and are shown to achieve explicit linear convergence rates.

%%% Our contribution.
\vspace{0.75ex}
\noindent\textbf{$\mathrm{(c)}$~Our contributions.}
To this end, the main contributions of this paper are summarized as follows.

\begin{compactitem}
\item[$\mathrm{(i)}$]
We propose a new parameter-free Halpern fixed-point algorithm for approximating a fixed point of \eqref{eq:FP}, where $T$ is a $\rho$-contractive mapping with $\rho \in (0,1)$. 
The proposed method does not require any prior knowledge of the contraction factor $\rho$. 
We establish explicit linear convergence rates both for the fixed-point residual $\norms{x^k - T(x^k)}$ and the distance to the fixed point $\norms{x^k - x^{\star}}$.

\item[$\mathrm{(ii)}$]
Inspired by \cite{he2024convergence}, we develop a new adaptive Halpern fixed-point algorithm for approximating a fixed point of \eqref{eq:FP}, where $T$ is a $\rho$-contractive mapping with $\rho \in (0,1]$. 
Unlike the first method, this algorithm requires only an upper bound of $\rho$ to provably guarantee a linear convergence and reduces to \cite[Algorithm~3.1]{he2024convergence} when $\rho=1$. 
We also establish explicit linear convergence rates for both the fixed-point residual $\norms{x^k - T(x^k)}$ and the distance to the fixed point $\norms{x^k - x^{\star}}$.

\item[$\mathrm{(iii)}$]
By combining the proposed fixed-point algorithms with Tikhonov regularization, we develop a new parameter-free method for solving the co-coercive equation \eqref{eq:CE} when $G$ is $\beta$-co-coercive. 
Building upon the convergence results established in (i) and (ii), we prove the convergence of the proposed method and derive the iteration complexity $\BigOs{\epsilon^{-1}\ln(\epsilon^{-1})}$ for computing an $\epsilon$-solution $x^k_{\epsilon}$ satisfying $\norms{G(x^k_{\epsilon})}\leq\epsilon$ for a given accuracy $\epsilon > 0$.

\item[$\mathrm{(iv)}$]
Exploiting the equivalence between Halpern's fixed-point iteration and Nesterov's accelerated schemes established in \cite{tran2022connection}, we derive a new parameter-free Nesterov-accelerated fixed-point algorithm for solving \eqref{eq:FP}. 
As a consequence of the results established in (i) and (ii), we prove that the proposed algorithm also enjoys a linear convergence rate whenever $\rho\in(0,1)$.
\end{compactitem}
%%%
\textbf{\textit{Discussion.}}
Let us further elaborate on the main contributions of this paper. 
%%
First, our work complements the recent developments in \cite{diakonikolas2020halpern,he2024convergence} by establishing explicit linear convergence rates for contractive mappings. 
This distinguishes our results from many existing works, where only asymptotic convergence is available without explicit convergence rates or iteration-complexity guarantees.
%%
Second, the proposed algorithms incur essentially the same per-iteration computational cost as the classical Halpern iteration \eqref{eq:HP_iteration_rule}, BP iterations, or KM methods. 
The additional computation consists only of a few vector operations, inner products, and norm evaluations, all of which require $\BigOs{p}$ operations.
%%
Third, although our theoretical analysis is developed for contractive mappings, the numerical experiments in Section~\ref{sec:num_examples} demonstrate that the proposed algorithms also perform remarkably well on nonexpansive mappings. 
In particular, they consistently outperform the existing methods considered in our experiments, including \cite{he2024convergence}.
%%
Fourth, by combining the proposed methods with Tikhonov regularization, we obtain new adaptive algorithms for solving co-coercive equations, or equivalently, fixed-point problems involving nonexpansive mappings. 
This substantially broadens the applicability of our approach to a variety of algorithms, including proximal-point, projection, and operator-splitting methods.
%%
Fifth, our Nesterov's accelerated variants further reveal an explicit equivalence between Halpern's fixed-point iteration and Nesterov's acceleration in the context of parameter-free and adaptive algorithms.
%%
Finally, although the numerical experiments reported in this paper are based on synthetic examples, they provide encouraging evidence that the proposed methods have the potential to be effective for a broad range of practical applications.

%%% (d) Related work.
\vspace{0.75ex}
\noindent\textbf{$\mathrm{(d)}$~Related work.}
As mentioned earlier, parameter-free and adaptive algorithms have been extensively studied in optimization and have recently gained significant attention due to their success in machine learning and artificial intelligence, see, e.g., \cite{defazio2022parameter,defazio2023learning,Duchi2011,KingmaB14,jordan2024muon,orabona2016coin}. 
In contrast, despite several recent advances, the development of provably convergent and practically efficient parameter-free algorithms for fixed-point problems and operator theory remains relatively limited.

Over the past decades, a number of researchers have investigated adaptive and parameter-free methods for fixed-point problems, as well as root-finding methods for monotone inclusions (MIs) and variational inequality problems (VIPs). 
%%%
In fixed-point theory, several adaptive iterative schemes have been proposed. 
For example, Colao and Marino \cite{colao2015krasnoselskii} studied an adaptive variant of the Krasnosel'ski\v{\i}-Mann (KM) iteration by introducing adaptive update rules for the relaxation parameters. 
Diakonikolas \cite{diakonikolas2020halpern} proposed a parameter-free Halpern fixed-point method for solving the co-coercive equation \eqref{eq:CE}, achieving a non-asymptotic convergence rate of $\BigOs{1/k}$. 
More recently, He et al.~\cite{he2024convergence} developed an adaptive Halpern fixed-point method for solving \eqref{eq:FP} with nonexpansive mappings. 
Their method eliminates the need for line search while demonstrating promising numerical performance on several representative examples. 
Building upon this idea, Lv et al.~\cite{lv2026preconditioned} proposed a preconditioned Halpern method with adaptive parameter updates.

For monotone inclusions and variational inequality problems, numerous adaptive algorithms based on line-search techniques have been developed, comprehensive accounts can be found in the monographs \cite{Facchinei2003,Konnov2001} and some recent works such as \cite{he2018totally,oyewole2022totally}. 
In particular, Malitsky and his collaborators proposed several adaptive extragradient-type methods equipped with line-search procedures for solving these problems, see, for example, \cite{malitsky2019golden,malitsky2020forward}. 
These algorithms have subsequently found applications in minimax optimization, game theory, and machine learning, see, e.g., \cite{daskalakis2018training,ioan2023relaxed,malitsky2016first}.

%%% (e) Paper organization.
\vspace{0.75ex}
\noindent\textbf{$\mathrm{(e)}$~Paper organization.}
The remainder of this paper is organized as follows.
Section~\ref{sec:prelim_results} reviews the background material and presents some preliminary results used throughout the paper.
Section~\ref{sec:linear_convergence_rate} introduces our first parameter-free Halpern fixed-point algorithm and establishes its linear convergence.
Section~\ref{sec:adaptive_Halpern_method} develops an adaptive Halpern fixed-point algorithm together with its convergence analysis.
Section~\ref{sec:NonExMapping} applies the proposed methods to solve co-coercive equations by combining them with Tikhonov regularization.
Section~\ref{sec:NesMethods} exploits the connection established in \cite{tran2022connection} to derive parameter-free Nesterov-accelerated variants of our methods.
Finally, Section~\ref{sec:num_examples} presents numerical experiments demonstrating the performance of the proposed algorithms and compares them with existing methods.


%%%%%%%%%%%%%%%%%%%%%%%%%%%%%%%%%%%%%%%%
%%%% 2. Preliminary results
%%%%%%%%%%%%%%%%%%%%%%%%%%%%%%%%%%%%%%%%
\section{Background and Preliminary Results}\label{sec:prelim_results}
We begin by reviewing the basic notation and concepts used throughout the paper, including fixed-point mappings and co-coercive equations. 
For further details, we refer to \cite{Bauschke2011,Facchinei2003}.
We then briefly recall two classical methods for approximating a fixed point of a $\rho$-contractive mapping.

%%% 2.1. Notations and basic concepts.
\subsection{Notations and basic concepts}\label{sec:notation}
Throughout this paper, $\Hil$ denotes a Hilbert space equipped with the standard inner product $\iprod{\cdot,\cdot}$ and the induced norm $\norms{\cdot}$.
For a function $f : \Hil \to\R$, $\nabla{f}$ denotes its gradient, $\partial{f}$ its subdifferential, and $\prox_{f}$ its proximal operator.
For a matrix $\mbf{Q}\in\R^{n\times p}$, we denote by $\norms{\mbf{Q}}_2$ its spectral norm.
Given two functions $g(t)$ and $h(t)$, we write $g(t) = \BigOs{h(t)}$ if there exist constants $M > 0$ and $t_0 \geq 0$ such that $g(t) \leq Mh(t)$ for all $t \geq t_0$.
We use $\widetilde{\mcal{O}}(g(t))$ to suppress polylogarithmic factors of $g(t)$.

Let $T : \Hil \to \Hil$ be a single-valued mapping.
We say that $T$ is $\rho$-Lipschitz continuous with Lipschitz constant $\rho > 0$ if $\norms{T(x) - T(y)} \leq \rho\norms{x - y}$ for all $x, y \in \Hil$.
If $\rho \in [0, 1)$, then $T$ is called $\rho$-contractive.
If $\rho = 1$, then $T$ is called nonexpansive.
Furthermore, if $\norms{T(x) - T(y)}^2 \leq \norms{x - y}^2 - \norms{x - T(x) - (y - T(y))}^2$ for all $x, y \in \Hil$, then $T$ is called firmly nonexpansive.
We denote by $\textrm{Fix}(T) := \sets{x^{\star} \in \Hil : x^{\star} = T(x^{\star})}$ the set of fixed points of $T$.
It is well known that if $T$ is $\rho$-contractive, then $\textrm{Fix}(T)$ is a singleton.

A single-valued mapping $G : \Hil \to \Hil$ is called $\beta$-co-coercive with co-coercivity parameter $\beta > 0$ if $\iprods{G(x) - G(y), x - y} \geq \beta\norms{G(x) - G(y)}^2$ for all $x, y \in \Hil$.
If $\beta = 0$, then $G$ is called monotone.
It is well known that if $G$ is $\beta$-co-coercive, then $T := \Id - 2\beta G$ is nonexpansive.
Conversely, if $T$ is nonexpansive, then $G := \Id - T$ is $\frac{1}{2}$-co-coercive.
Moreover, every $\beta$-co-coercive mapping is also $\frac{1}{\beta}$-Lipschitz continuous.
We denote by $\mathrm{zer}(G) := \sets{x^{\star} \in \Hil : G(x^{\star}) = 0}$ the solution set of  $G(x)=0$.

%%%% 2.2. Halpern fixed-point iteration and tight convergence rates
\subsection{Halpern fixed-point iteration and tight convergence rates}\label{sec:HP_convergence}
Let $T : \Hil \to \Hil$ be a $\rho$-contractive mapping (also referred to as a $\rho$-Lipschitz continuous mapping) for some $\rho \in (0,1]$.
The Halpern fixed-point iteration for approximating a fixed point $x^{\star} \in \mathrm{Fix}(T)$ of $T$, i.e., $x^{\star} = T(x^{\star})$, is defined as follows.
Choose an anchor point $u \in \Hil$ and an initial point $x^0 \in \Hil$ (where $x^0$ may differ from $u$).
Then, for each iteration $k \geq 0$, update
\begin{equation}\label{eq:HP_iteration}
x^{k+1} := \lambda_k u + (1 - \lambda_k)T(x^k),
\tag{HP}
\end{equation}
where, following \cite{park2022exact}, the parameter $\lambda_k$ is chosen as
$\lambda_k = \frac{1}{1 + \varphi_k}$ with $\varphi_k := \sum_{i=1}^k\rho^{-2i}$.
For simplicity, throughout the remainder of the paper we set the anchor point to be the initial point, i.e., $u = x^0$.

The following theorem summarizes the convergence of the iteration \eqref{eq:HP_iteration}; see \cite[Corollary~3.3]{park2022exact}.

%%% Theorem 2.1.
\begin{theorem}[\cite{park2022exact}]\label{th:linear_convergence_of_HP}
Let $T : \Hil \to \Hil$ be a $\rho$-contractive mapping for some $\rho \in (0, 1]$, and let $x^{\star} \in \mathrm{Fix}(T)$ be a fixed point of $T$.
Let $\sets{x^k}$ be the sequence generated by \eqref{eq:HP_iteration} with $x^0 := u$.
Then
\begin{equation}\label{eq:th11_convergence_bound}
\norms{x^k - T(x^k)} \leq \frac{(1+\rho)}{\sum_{i=0}^{k-1}\rho^{-i} }\norms{x^0 - x^{\star}} =  \frac{(1-\rho^2)\rho^{k-1}\norms{x^0 - x^{\star}}}{1 - \rho^k} \leq (1+\rho)\rho^{k-1}\norms{x^0 - x^{\star}}, \quad \forall k \geq 1.
\end{equation}
\end{theorem}
%%%

When $\rho \in (0,1)$, i.e., $T$ is contractive, Theorem~\ref{th:linear_convergence_of_HP} establishes a linear convergence rate for the fixed-point residual $\norms{x^k - T(x^k)}$. On the other hand, when $\rho = 1$, i.e., $T$ is nonexpansive, the convergence bound \eqref{eq:th11_convergence_bound} reduces to
\begin{equation*}
\norms{x^k - T(x^k)} \leq \frac{2\norms{x^0 - x^{\star}}}{k+1},
\end{equation*}
which yields the sublinear convergence rate $\BigOs{1/k}$. 
As shown in \cite{park2022exact}, the rate stated in Theorem~\ref{th:linear_convergence_of_HP} is optimal and cannot be improved in the worst case.
If $T$ is $\rho$-contractive for $\rho \in (0, 1)$, then, by the relation $\norms{x^k - x^{\star}} \leq \frac{1}{1-\rho}\norms{x^k - T(x^k)}$ (see the proof in Theorem~\ref{th:convergence_theorem1}), we also get $\norms{x^k - x^{\star}} \leq \frac{(1+\rho)\rho^{k-1}}{1-\rho}\norms{x^0 - x^{\star}}$.
Therefore, $\set{x^k}$ also strongly converges to the unique fixed point $x^{\star}$ of $T$.

%%%%%%%%%%%%%%%%%%%%%%%%%%%%%%%%%%%%%%%%
%%%% 3. Linear convergence of Halpern's fixed-point iteration.
%%%%%%%%%%%%%%%%%%%%%%%%%%%%%%%%%%%%%%%%
\section{Main Result 1: Parameter-Free Halpern Fixed-Point Method}\label{sec:linear_convergence_rate}
In this section, we build upon the result of Theorem~\ref{th:linear_convergence_of_HP} to develop a parameter-free Halpern fixed-point method for approximating a fixed point of a $\rho$-contractive mapping $T$. 
We first derive the proposed method and present its implementation in Algorithm~\ref{alg:adaptive_HP0}. We then establish its linear convergence.

%%% Here is our proposed stepsize.
\subsection{The derivation of the algorithm}\label{subsec:derivation_of_A1}
Our main goal is to develop a parameter-free variant of \eqref{eq:HP_iteration} for the contractive case.
Suggested by the update rule $\varphi_k := \sum_{i=1}^k\rho^{-2i}$  of $\varphi_k$ in \eqref{eq:HP_iteration}, by the $\rho$-contractivity of $T$, it is natural to replace $\varphi_k$ by the following upper bound (when it is well-defined):
\begin{equation}\label{eq:new_stepsize0}
\bar{\varphi}_k := \sum_{i=1}^k\frac{\norms{x^i - x^{i-1}}^{2i}}{\norms{T(x^i) - T(x^{i-1})}^{2i}}.
\end{equation}
Indeed, by the $\rho$-contractivity  of $T$, we have $\norms{T(x^i) - T(x^{i-1})} \leq \rho\norms{x^i - x^{i-1}}$ for all $i = 1,\cdots, k$.
If $\norms{T(x^i) - T(x^{i-1})} > 0$ for all $i=1,\cdots, k$, then $\frac{\norms{x^i - x^{i-1}}^{2i}}{\norms{T(x^i) - T(x^{i-1})}^{2i}} \geq \frac{1}{\rho^{2i}}$ for $i=1,\cdots, k$.
Therefore, we can claim that
\begin{equation}\label{eq:new_stepsize0_bound}
\bar{\varphi}_k \geq \sum_{i=1}^k\rho^{-2i} =: \varphi_k.
\end{equation}
However, we may encounter $\norms{T(x^i) - T(x^{i-1})} = 0$ for some $i \in \sets{1, \cdots, k}$.
To avoid this situation, we first notice that $\varphi_k = \varphi_{k-1} + \rho^{-2k}$ for all $k \geq 0$ and, by convention, $\varphi_{-1} = 0$.
Next, we can replace $\varphi_k$ by $\bar{\varphi}_k$ updated as follows:
\begin{equation}\label{eq:new_stepsize}
\bar{\varphi}_k := \bar{\varphi}_{k-1} + \omega_k^{2k}, \quad \textrm{where} \quad \omega_k \geq \frac{1}{\rho} \quad \textrm{for all $k \geq 0$ and $\bar{\varphi}_{-1} := 0$}.
\end{equation}
Clearly, $\bar{\varphi}_k$ is also an upper bound of $\varphi_k$.

Now, we first choose $\omega_0$ such that $\omega_0 \geq \frac{1}{\rho}$.
By the $\rho$-contractivity of $T$, one way to choose the initial value $\omega_0$ is as follows:
\begin{equation}\label{eq:omega_0_choice}
\omega_0 := \frac{1}{\rho_0} \geq \frac{1}{\rho},\quad \textrm{where} \quad \rho_0 := \frac{\norms{T(x^0) - T(x^0 + v)}}{\norms{v}} \in (0,  \rho].
\end{equation}
for a given $v \neq 0$ in $\Hil$ such that $T(x^0 + v) \neq T(x^0)$.

For $k \geq 0$, given $\omega_k$, we can update $\omega_k$ as follows:
\begin{equation}\label{eq:update_omega_k}
\omega_{k+1} := \begin{cases}
\max\set{ \omega_k, \frac{\norms{x^{k+1} - x^k } }{ \norms{T(x^{k+1}) - T(x^{k})} } } & \textrm{if $\frac{\norms{x^{k+1} - x^k } }{ \norms{T(x^{k+1}) - T(x^{k})} } \leq \bar{\omega}$}, \\
\omega_k & \textrm{otherwise}.
\end{cases}
\end{equation}
Here, we set an upper bound $\bar{\omega} \in (0, +\infty)$ on $\frac{ \norms{x^{k+1} - x^k } }{ \norms{T(x^{k+1}) - T(x^{k})} }$ to prevent the explosion of this ratio.
We can heuristically choose $\bar{\omega} := \frac{\Lambda}{\rho_0} > \frac{1}{\rho}$ for some $\Lambda > 1$ sufficiently large, e.g., $\bar{\omega} := \frac{2^4}{\rho_0}$.


Since $\omega_0 \geq \frac{1}{\rho}$ and $\frac{\norms{x^{k+1} - x^k } }{ \norms{T(x^{k+1}) - T(x^{k})} } \geq \frac{1}{\rho}$, it is obvious to see that $\omega_k \geq \frac{1}{\rho}$ for all $k \geq 0$.
Using $\bar{\varphi}_k$ from \eqref{eq:update_omega_k} to replace $\varphi_k$ in the scheme \eqref{eq:HP_iteration}, we obtain a parameter-free variant of this algorithm.

Note that we can also replace the update rule \eqref{eq:update_omega_k} by the following one:
\begin{equation}\label{eq:update_omega_k_modified}
\omega_{k+1} := \begin{cases}
\min\set{ \omega_k, \frac{\norms{x^{k+1} - x^k } }{ \norms{T(x^{k+1}) - T(x^{k})} } } & \textrm{if $\frac{\norms{x^{k+1} - x^k } }{ \norms{T(x^{k+1}) - T(x^{k})} } \leq \bar{\omega}$}, \\
\omega_k & \textrm{otherwise}.
\end{cases}
\end{equation}
Then, we have the following fact, $\omega_k \geq \frac{1}{\rho}$ and $\omega_{k+1} \leq \omega_k$ for all $k \geq 0$.
The relation $\omega_{k+1} \leq \omega_k$ for all $k \geq 0$ is a direct consequence of \eqref{eq:update_omega_k_modified}.
We prove $\omega_k \geq \frac{1}{\rho}$ by induction. 
For $k = 0$, we have $\omega_0 := \frac{1}{\rho_0} \geq \frac{1}{\rho}$.
Assume that  $\omega_k \geq \frac{1}{\rho}$ for $k \geq 1$.
We prove that $\omega_{k+1} \geq \frac{1}{\rho}$.
Indeed, from \eqref{eq:update_omega_k_modified} we have two possibilities: (i) $\omega_{k+1} = \omega_k \geq \frac{1}{\rho}$ or (ii) $\omega_{k+1} = \frac{\norms{x^{k+1} - x^k } }{ \norms{T(x^{k+1}) - T(x^{k})} } \geq \frac{1}{\rho}$.
Therefore, we obtain our conclusion that $\omega_k \geq \frac{1}{\rho}$ for all $k \geq 0$. 

Now, we combine all the steps discussed above algorithmically to obtain Algorithm~\ref{alg:adaptive_HP}.

%%%% Algorithm 1.
\begin{algorithm}[!htbp]
\caption{(Parameter-Free Halpern Fixed-Point Algorithm)}\label{alg:adaptive_HP0}
\begin{algorithmic}[1]
\State\textbf{Initialization:} Choose an initial point $x^0\in\mathcal H$.
Estimate $\omega_0$ as in \eqref{eq:omega_0_choice} and set $\varphi_{-1} := 0$.
\State \textbf{For $k = 0, 1, 2, \cdots$, perform}
\State\hspace{3ex}\label{alg1:step3}If ~$\norms{r^k} = 0$ for $r^k := x^k - T(x^k)$,  then TERMINATE.
\State\hspace{3ex}Update $\bar{\varphi}_{k} := \bar{\varphi}_{k-1} + \omega_k^{2k}$.
\State\hspace{3ex}\label{alg1:step5}Update
\begin{equation*}
x^{k+1} :=  \frac{1}{ \bar{\varphi}_k + 1} x^0 + \frac{ \bar{\varphi}_k}{\bar{\varphi}_k + 1}T(x^k).
\end{equation*}
\State\hspace{3ex}\label{alg1:step6}Update $\omega_{k+1}$ as in \eqref{eq:update_omega_k} or \eqref{eq:update_omega_k_modified}.
\State\textbf{End For}
\end{algorithmic}
\end{algorithm}

%%%
\noindent
We make the following remarks regarding Algorithm~\ref{alg:adaptive_HP0}.
%%
\begin{compactitem}
\item Algorithm~\ref{alg:adaptive_HP0} is completely parameter-free, as it does not require any prior knowledge of the contraction factor $\rho$ of $T$. 
Moreover, it is straightforward to implement since it does not involve any inner loop, line-search, or bisection procedure.

\item The stopping criterion in Step~\ref{alg1:step3} is primarily intended for theoretical analysis. 
In practice, given a prescribed accuracy $\epsilon > 0$, we terminate Algorithm~\ref{alg:adaptive_HP0} whenever $\norms{r^k} \leq \epsilon$ or $\norms{r^k} \leq \epsilon\cdot\max\set{1, \norms{r^0}}$.
As a safeguard, we can also set the maximum number of iterations to $k_{\max}$.

\item The update rule for $\omega_{k+1}$ in Step~\ref{alg1:step6} requires only the evaluation of the two norms $\norms{x^{k+1} - x^k}$ and $\norms{T(x^{k+1}) - T(x^k)}$, together with two elementary arithmetic operations. 
Consequently, it introduces only negligible computational overhead compared with the classical Halpern iteration.
\end{compactitem}
%%
Overall, the per-iteration computational complexity of Algorithm~\ref{alg:adaptive_HP0} is essentially the same as that of the classical Halpern, Banach-Picard (BP), and Krasnosel'ski\v{\i}-Mann (KM) fixed-point iterations.

%%% 3.2. Convergence analysis.
\subsection{Convergence analysis}\label{subsec:convergence_guarantee_of_A1}
The following theorem proves linear convergence rates of Algorithm~\ref{alg:adaptive_HP0}.

%%% Theorem 1.
\begin{theorem}\label{th:convergence_theorem1}
Let $T : \Hil \to \Hil$ be $\rho$-contractive for some $\rho \in (0, 1)$ and $x^{\star}$ be the unique fixed-point of $T$.
Let $\sets{x^k}$ be generated by Algorithm~\ref{alg:adaptive_HP0}.
Then, we have
\begin{equation}\label{eq:convergence_bound1}
\norms{x^k - T(x^k)} \leq  C_{\rho} \rho^k \norms{x^0 - T(x^0)} \quad \textrm{and} \quad \norms{x^k - x^{\star}} \leq \frac{C_{\rho}(1+\rho)}{1-\rho} \rho^k \norms{x^0 - x^{\star}},
\end{equation}
where $C_{\rho} := 1 + \frac{1+\rho}{(1-\rho)^2\rho}$.
Consequently, $\sets{x^k}$ strongly converges to $x^{\star}$ with a linear rate.
\end{theorem}

%%% Proof of Theorem 1.
\begin{proof}
Denote by $\lambda_k := \frac{1}{\varphi_k + 1}$, $r^k := x^k - T(x^k)$, and $R_k := \norms{x^k - x^0}$.
First, since $\bar{\varphi}_k \geq 0$ for all $k \geq 0$, we have $1 + \bar{\varphi}_k = 1 + \bar{\varphi}_{k-1} + \omega^{2k} \geq \omega^{2k}$ for all $k \geq 1$.
Second, by the choice of $\rho_0$ in \eqref{eq:omega_0_choice}, we have $\omega_0 \geq \frac{1}{\rho}$.
Third, if $\omega_k$ is updated by \eqref{eq:update_omega_k}, then we have 
\begin{equation*}
\omega_{k+1} \geq \omega_k \quad \textrm{and} \quad \omega_{k+1} \geq \frac{\norms{x^{k+1} - x^k } }{ \norms{T(x^{k+1}) - T(x^{k})} } \geq \frac{1}{\rho}.
\end{equation*}
Alternatively, if $\omega_k$ is updated by \eqref{eq:update_omega_k_modified}, then as we have shown earlier $\omega_k \geq \frac{1}{\rho}$.
In both cases, by induction, we have $\omega_k \geq \frac{1}{\rho}$ for all $k \geq 0$.
Thus, we conclude that $1 + \bar{\varphi}_k \geq \rho^{-2k}$, leading to $\lambda_k := \frac{1}{1+ \bar{\varphi}_k} \leq \rho^{2k}$ for all $k \geq 0$.

Now, from Step \ref{alg1:step5} of Algorithm~\ref{alg:adaptive_HP0} and the definition of $\lambda_k := \frac{1}{\bar{\varphi}_k + 1} \in (0, 1)$, we have $x^{k+1} = \lambda_kx^0 + (1-\lambda_k)T(x^k)$, leading to $x^{k+1} - x^0 = (1-\lambda_k)(T(x^k) - x^0)$.
Using the triangle inequality, this relation, and the $\rho$-contractivity of $T$, we can show that
\begin{equation*}
\arraycolsep=0.2em
\begin{array}{lcl}
R_{k+1} &:= & \norms{x^{k+1} - x^0} = (1-\lambda_k)\norms{T(x^k) - x^0} \leq \norms{T(x^k) - x^0} \vspace{1ex}\\
& \leq & \norms{T(x^k) - T(x^0)} + \norms{x^0 - T(x^0)} \vspace{1ex}\\
& \leq & \rho\norms{x^k - x^0} + \norms{r^0} \vspace{1ex}\\
& = & \rho R_k + \norms{r^0}.
\end{array}
\end{equation*}
By induction, we obtain $R_k \leq \rho^k R_0 + \norms{r^0}\sum_{i=0}^{k-1}\rho^i = \frac{1-\rho^k}{1-\rho}\norms{r^0} \leq \frac{\norms{r^0}}{1-\rho}$  since $R_0 = 0$.
Using again the $\rho$-contractivity of $T$ and this inequality, we can show that
\begin{equation}\label{eq:th21_proof1}
\arraycolsep=0.2em
\begin{array}{lcl}
\norms{x^0 - T(x^k)} & \leq & \norms{x^0 - T(x^0)} + \norm{T(x^0) - T(x^k)} \leq \norms{r^0} + \rho R_k \vspace{1ex}\\
& \leq & \Big(1 + \frac{\rho}{1-\rho}\Big)\norms{r^0} = \frac{1}{1-\rho}\norms{r^0}.
\end{array}
\end{equation}
Next, using one more time the $\rho$-contractivity of $T$ and the triangle inequality, we can derive that
\begin{equation}\label{eq:th21_proof2} 
\arraycolsep=0.2em
\begin{array}{lcl}
\norms{r^{k+1}} & = & \norms{x^{k+1} - T(x^{k+1})} = \norms{x^{k+1} - T(x^k) + [T(x^k) - T(T(x^k))] + [T(T(x^k)) - T(x^{k+1})]} \vspace{1ex}\\
& \leq & \norms{T(x^k) - T(T(x^k))} + \norms{x^{k+1} - T(x^k)} +  \norms{T(T(x^k)) - T(x^{k+1})} \vspace{1ex}\\
& \leq & \rho\norms{r^k} +  (1 + \rho)\norms{x^{k+1} - T(x^k)}.
\end{array}
\end{equation}
On the other hand, again, by Step \ref{alg1:step5} of Algorithm~\ref{alg:adaptive_HP0}, we have 
\begin{equation*} 
\arraycolsep=0.2em
\begin{array}{lcl}
\norms{x^{k+1} - T(x^k)} &= & \norms{ \lambda_kx^0 + (1-\lambda_k)T(x^k) - T(x^k)} = \norms{\lambda_k(x^0 - T(x^k))} = \lambda_k\norms{x^0 - T(x^k)}.
\end{array}
\end{equation*}
Combining this expression, \eqref{eq:th21_proof2}, and \eqref{eq:th21_proof1}, we get 
\begin{equation*} 
\arraycolsep=0.2em
\begin{array}{lcl}
\norms{r^{k+1}} \leq \rho\norms{r^k} + (1+\rho)\lambda_k\norms{x^0 - T(x^k)} \leq \rho\norms{r^k} + \frac{1+\rho}{1-\rho}\lambda_k\norms{r^0}.
\end{array}
\end{equation*}
As we have shown above that $\lambda_k \leq \rho^{2k}$, the last inequality leads to
\begin{equation}\label{eq:th21_proof3} 
\arraycolsep=0.2em
\begin{array}{lcl}
\norms{r^{k+1}} \leq \rho\norms{r^k} + \frac{(1+\rho)\norms{r^0}}{1-\rho}   \rho^{2k} = \rho \norms{r^k} + C_0\norms{r^0} \rho^{2k},
\end{array}
\end{equation}
where $C_0 : = \frac{1+\rho}{1-\rho}$.
By induction, we can show from \eqref{eq:th21_proof3} that
\begin{equation*} 
\arraycolsep=0.2em
\begin{array}{lcl}
\norms{r^k} & \leq & \rho^k\norms{r^0} + C_0\norms{r^0} \sum_{i=0}^{k-1}\rho^{k-1-i}\rho^{2i} = \rho^k\norms{r^0} + C_0\norms{r^0} \rho^{k-1}\sum_{i=0}^{k-1}\rho^i \vspace{1ex}\\
& \leq & \Big( \rho + \frac{1+\rho}{(1-\rho)^2}\Big)\rho^{k-1} \norms{r^0} = \Big(1 + \frac{1+\rho}{\rho(1-\rho)^2}\Big)\rho^{k} \norms{r^0}.
\end{array}
\end{equation*}
This is exactly the first bound of \eqref{eq:convergence_bound1} by using $r^k := x^k - T(x^k)$.

Finally, since $T$ is $\rho$-contractive, we have  
 \begin{equation*}
     \norms{ x^k - x^{\star} }  = \norms{ x^k - T(x^k) + T(x^k) - T(x^{\star}) } \leq \norms{ x^k - T(x^k) } + \rho \norms{ x^k - x^{\star} }.
 \end{equation*}
 Thus, we get $\norms{ x^k - x^{\star} }  \leq \frac{1}{1-\rho} \norms{ x^k - T(x^k)}$. 
 Moreover, by $x^{\star} = T(x^{\star})$ and  the $\rho$-contractivity of $T$, we can easily show that $\norms{r^0} = \norms{x^0 - T(x^0)} \leq \norms{x^0 - x^{\star}} + \norms{T(x^0) - T(x^{\star}} \leq (1+\rho)\norms{x^0 - x^{\star}}$.
Combining the last two inequalities and the first bound of  \eqref{eq:convergence_bound1}, we get the second bound of  \eqref{eq:convergence_bound1}.
Since $\rho^k \to 0$ as $k \to \infty$, the second bound of  \eqref{eq:convergence_bound1} shows that $\sets{x^k}$ converges strongly to $x^{\star} \in \mathrm{Fix}(T)$.
\end{proof}
%%% End of Proof.

%%%%%%%%%%%%%%%%%%%%%%%%%%%%%%%%%%%%%%%%%%%%%%%%
%%% 4. Adaptive Halpern Fixed-Point Method
%%%%%%%%%%%%%%%%%%%%%%%%%%%%%%%%%%%%%%%%%%%%%%%%
\section{Main Result 2: Adaptive Halpern Fixed-Point Method}\label{sec:adaptive_Halpern_method}
Inspired by \cite{he2024convergence}, we develop in this section an adaptive Halpern fixed-point method for approximating the unique fixed point of a $\rho$-contractive mapping. 
Unlike the nonexpansive setting considered in \cite{he2024convergence}, where no problem-dependent parameter is required, our framework naturally involves the contraction parameter $\rho$. 
Our objective is to minimize the reliance on this parameter as much as possible. 
In particular, the proposed method requires only an upper bound of $\rho$, rather than its exact value. 
Although the resulting algorithm is not theoretically completely parameter-free, as is Algorithm~\ref{alg:adaptive_HP0}, it significantly reduces the amount of problem-specific information needed in practice. 
Following the terminology introduced in \cite{he2024convergence}, we refer to this approach as an \emph{adaptive Halpern fixed-point method}, since it adaptively updates the anchor parameter throughout the iterations.

%%% 4.1. The derivation of the algorithm
\subsection{The derivation of the algorithm}\label{subsec:A2_derivation}
Following the same idea as in \cite{he2024convergence}, we first define
\begin{equation}\label{eq:residuals}	
	r^k := x^k-T(x^k) \quad \text{and} \quad s^k := x^0-T(x^k).
\end{equation}	
Next, starting from the Halpern fixed-point iteration \eqref{eq:HP_iteration} with $\lambda_k = \hat{\lambda}_k := \frac{1}{\hat{\varphi}_k + 1}$ for a given $\hat{\varphi}_k > 0$ and $u = x^0$, we can rewrite it as
\begin{equation}\label{eq:HP_iteration2}
x^{k+1} := \hat{\lambda}_k x^0 + (1 - \hat{\lambda}_k)T(x^k) = \frac{1}{\hat{\varphi}_k + 1}x^0 + \frac{\hat{\varphi}_k}{\hat{\varphi}_k + 1}T(x^k).
\end{equation}
Using this formula and \eqref{eq:residuals}, we can show that
\begin{equation*}
\arraycolsep=0.2em
\begin{array}{lcl}
x^{k+1} - T(x^k) & = & \hat{\lambda}_kx^0 + (1-\hat{\lambda}_k)T(x^k) - T(x^k) = \hat{\lambda}_k(x^0 - T(x^k)) = \hat{\lambda}_ks^k, \vspace{1ex}\\
x^{k+1} - x^k & = & \hat{\lambda}_k x^0 +  (1-\hat{\lambda}_k)T(x^k) - x^k = \hat{\lambda}_k(x^0 - T(x^k))  - (x^k - T(x^k)) = \hat{\lambda}_ks^k - r^k.
\end{array}
\end{equation*}
Let $\bar{\rho} \in (0, 1)$ be an upper bound of $\rho$, i.e., $0 < \rho \leq \bar{\rho} < 1$.
By the $\rho$-contractivity of $T$, we have $\norms{T(x^{k+1}) - T(x^k)} \leq \rho\norms{x^{k+1} - x^k} \leq \bar{\rho}\norms{x^{k+1} - x^k}$ for any $k \geq 0$.
Then, utilizing the above two expressions, and the last inequality, we can show that 
\begin{equation*}
\arraycolsep=0.2em
\begin{array}{lcl}
0 &\leq & \bar{\rho}^2\norms{x^{k+1} - x^k}^2 - \norms{T(x^{k+1}) - T(x^k)}^2 \vspace{1ex}\\
& = & \bar{\rho}^2\norms{\hat{\lambda}_ks^k - r^k}^2 - \norms{x^{k+1} - T(x^{k+1}) - (x^{k+1} - T(x^k))}^2 \vspace{1ex}\\
& = & \bar{\rho}^2\norms{\hat{\lambda}_ks^k - r^k}^2 - \norms{r^{k+1} - \hat{\lambda}_ks^k}^2 \vspace{1ex}\\
& = & \bar{\rho}^2\norms{\hat{\lambda}_ks^k - r^k}^2 - \norms{r^{k+1}}^2 + 2\hat{\lambda}_k\iprods{s^k, r^{k+1}} - \hat{\lambda}_k^2\norms{s^k}^2.
\end{array}
\end{equation*}
This inequality leads to 
\begin{equation*}
\arraycolsep=0.2em
\begin{array}{lcl}
\norms{r^{k+1}}^2 &\leq & 2\hat{\lambda}_k\iprods{s^k,  r^{k+1}} + \bar{\rho}^2\norms{\hat{\lambda}_ks^k - r^k}^2  - \hat{\lambda}_k^2\norms{s^k}^2.
\end{array}
\end{equation*}
Therefore, if we impose 
\begin{equation}\label{eq:descent_cond}
\bar{\rho} \norms{\hat{\lambda}_ks^k - r^k}  \leq \hat{\lambda}_k\norms{s^k},
\end{equation}
then the last inequality reduces to
\begin{equation}\label{eq:descent_pro}
\norms{r^{k+1} }^2 \leq  2\hat{\lambda}_k\iprods{s^k, r^{k+1}}.
\end{equation}
Our objective is to develop an adaptive update rule for $\hat{\lambda}_k$ based on the descent condition \eqref{eq:descent_cond}. 
However, this condition still depends on the upper bound $\bar{\rho}$ of the contraction factor. 
To eliminate this dependence, we first establish a sequence of technical lemmas that will lead to the desired update rule for $\hat{\lambda}_k$. 
We begin with the following lemma.

%%% Lemma 1.
\begin{lemma}\label{le:key_properties}
Let $\sets{x^k}_{k \geq 0}$ be updated by \eqref{eq:HP_iteration2} and $\sets{\hat{\varphi}_k} \subset (0, +\infty)$ be such that $\hat{\lambda}_k = \frac{1}{1 + \hat{\varphi}_k}$.
Let $r^k$ and $s^k$ be defined by \eqref{eq:residuals}.
Assume that $\bar{\rho} \norms{\hat{\lambda}_ks^k - r^k}  \leq \hat{\lambda}_k\norms{s^k}$ for all $k \geq 0$ as in \eqref{eq:descent_cond}.
Then
\begin{itemize}
\item[{\rm(i)}]  $\bar{\rho}\norms{x^{k+1} -  x^k} \leq \norms{ x^{k+1} - T(x^k) }$.
\item[{\rm(ii)}] $\norms{ r^{k+1} }^2 \leq 2 \hat{\lambda}_k \iprods{s^k, r^{k+1}}$.
\item[{\rm(iii)}] $\norms{ r^{k+1}}^2 \leq \frac{2}{\hat{\varphi}_k} \iprods{ r^{k+1}, x^0 - x^{k+1}}$.
\end{itemize}
\end{lemma}

%%% Proof of Lemma 1.
\begin{proof}
(i). From \eqref{eq:residuals} and \eqref{eq:HP_iteration2}, we get $x^{k+1} = T(x^k) + \hat{\lambda}_k s^k$, which immediately implies that $\norms{ x^{k+1} - T(x^k)} = \hat{\lambda}_k \norms{s^k}$ and $x^k - x^{k+1} = (x^k - T(x^k)) - \hat{\lambda}_k s^k = r^k - \hat{\lambda}_k s^k$.
Substituting these two expressions into $\bar{\rho} \norms{\hat{\lambda}_ks^k - r^k}  \leq \hat{\lambda}_k\norms{s^k}$, we obtain $\bar{\rho} \norms{ x^k - x^{k+1}} \leq \norms{ x^{k+1} - T(x^k)}$, which proves (i).
			
(ii) Next, utilizing $r^k$ and $s^k$ from \eqref{eq:residuals}, and \eqref{eq:HP_iteration2}, we can show that
\begin{equation*}
r^{k+1} = x^{k+1} - T(x^{k+1}) = T(x^k) - T(x^{k+1}) + \hat{\lambda}_k s^k,
\end{equation*}
which is equivalent to $r^{k+1} - \hat{\lambda}_k s^k = T(x^k) - T(x^{k+1})$.

Now, since $\norms{T(x^{k+1}) - T(x^k)} \leq \rho\norms{x^{k+1} - x^k} \leq \bar{\rho}\norms{x^{k+1} - x^k}$, using (i), we can deduce that
\begin{equation*}
\norms{ r^{k+1} - \hat{\lambda}_k s^k } = \norms{ T(x^{k+1}) - T(x^k)} \leq \bar{\rho} \norms{ x^{k+1} - x^k } \overset{\tiny\mathrm{(i)}}{\leq} \norms{ x^{k+1} - T(x^k) } =  \hat{\lambda}_k  \norms{ s^k }.
\end{equation*}
Squaring both sides and cancelling the common term $\hat{\lambda}_k^2\norms{s^k}^2$, we get (ii).
			
(iii). From \eqref{eq:HP_iteration2}, we have $x^0 - x^{k+1} = (1 - \hat{\lambda}_k) (x^0 - T(x^k)) = (1 - \hat{\lambda}_k)s^k$.
Combining this relation and  (ii), we can deduce that
\begin{equation*}
\norms{ r^{k+1} }^2 \leq \frac{2\hat{\lambda}_k}{1 - \hat{\lambda}_k} \iprods{ r^{k+1}, x^0 - x^{k+1}}.
\end{equation*}
Since $\hat{\lambda}_k = \frac{1}{\hat{\varphi}_k + 1}$, we obtain (iii) from the last inequality.
\end{proof}
%%% End of Proof.

Next, we establish a lower bound on $\hat{\varphi}_k$ to prove a linear convergence rate of our method.

%%% Lemma 2.
\begin{lemma}\label{le:varphi_k_bounds}
Let $\sets{x^k}_{k \geq 0}$ be updated by \eqref{eq:HP_iteration2} and $\sets{\hat{\varphi}_k} \subset (0, +\infty)$ be such that $\hat{\lambda}_k = \frac{1}{1 + \hat{\varphi}_k}$.
Let $r^k$ and $s^k$ be defined by \eqref{eq:residuals}.
Then, the following statements hold.
\begin{itemize}
\item[{\rm(i)}] If $\bar{\rho} \norms{\hat{\lambda}_ks^k - r^k}  = \hat{\lambda}_k\norms{s^k}$ holds for $k \geq 0$, then we have $\hat{\varphi}_k \geq \left(\frac{1 + \bar{\rho}}{2\bar{\rho}}\right)^k\left(\hat{\varphi}_0 + \frac{2}{1 - \bar{\rho}} \right) - \frac{2}{1 - \bar{\rho}}$.

\item[{\rm(ii)}] If $\bar{\rho} \norms{\hat{\lambda}_ks^k - r^k}  \leq \hat{\lambda}_k\norms{s^k}$ holds for all $k \geq 0$ and $2\omega\hat{\lambda}_k\iprods{s^k, r^k} \leq \norms{r^k}^2$ hold for all $k \geq 1$ and a fixed $\omega \in (1, +\infty)$, then we have $\hat{\varphi}_k \geq \omega^k \big(\hat{\varphi}_0 +  \frac{2\omega-1}{\omega - 1}\big) -  \frac{2\omega-1}{\omega - 1}$ for all $k \geq 0$.
\end{itemize}
\end{lemma}

%%% Proof of Lemma 2.
\begin{proof}
(i). First, utilizing $s^{k+1}$ and $r^{k+1}$ from \eqref{eq:residuals}, and \eqref{eq:HP_iteration2}, we can easily show that
\begin{equation*}
s^{k+1} = x^0 - T(x^{k+1})  = s^k + T(x^k) - T(x^{k+1})  = r^{k+1} + (1 - \hat{\lambda}_k)s^k.
\end{equation*}	
This relation implies that $\iprods{ r^{k+1}, s^{k+1} } = \norms{ r^{k+1} }^2 + (1 - \hat{\lambda}_k)\iprods{ r^{k+1}, s^k}$.

Next, applying Lemma~\ref{le:key_properties}(ii), we get
\begin{equation*}
\iprods{ r^{k+1}, s^{k+1} }  \geq \norms{ r^{k+1} }^2 + \frac{1 - \hat{\lambda}_k}{ 2 \hat{\lambda}_k} \norms{ r^{k+1} }^2 = \frac{1 + \hat{\lambda}_k}{2\hat{\lambda}_k} \norms{ r^{k+1} }^2.
\end{equation*}
For all $k \geq 0$, let us define $U_k := \frac{\iprods{r^k, s^k}}{\norms{r^k}^2}$.
Then, we have $U_{k+1} = \frac{ \iprods{ r^{k+1}, s^{k+1} }}{ \norms{ r^{k+1} }^2 }$ and the last inequality implies that
\begin{equation}\label{eq:B_k}
U_{k+1}	\ge \frac{1 + \hat{\lambda}_k}{2\hat{\lambda}_k}.
\end{equation}
Since $\bar{\rho} \norms{\hat{\lambda}_ks^k - r^k}  = \hat{\lambda}_k\norms{s^k}$ holds for all $k \geq 0$, we have $\bar{\rho}^2\norms{\hat{\lambda}_{k+1}s^{k+1} - r^{k+1} }^2 = \hat{\lambda}_{k+1}^2\norms{s^{k+1}}^2$, which is equivalent to
\begin{equation}\label{eq:lm1_proof1}
2\hat{\lambda}_{k+1}\frac{\iprods{ r^{k+1}, s^{k+1} } }{\norms{ r^{k+1} }^2}  + \frac{(1 - \bar{\rho}^2)}{ \bar{\rho}^2} \hat{\lambda}_{k+1}^2\frac{\norms{ s^{k+1} }^2}{\norms{ r^{k+1}}^2 } = 1.
\end{equation}
By the Cauchy-Schwarz inequality and the definition of $U_k$, we can show that
\begin{equation}\label{eq:lm1_proof2}
U^2_{k+1} = \frac{\iprods{ r^{k+1}, s^{k+1}}^2}{\norms{ r^{k+1} }^4} \leq \frac{\norms{ r^{k+1}}^2 \norms{ s^{k+1}}^2}{\norms{ r^{k+1}}^4} = \frac{\norms{ s^{k+1}}^2}{\norms{ r^{k+1}}^2}.
\end{equation}
Combining \eqref{eq:lm1_proof1} and \eqref{eq:lm1_proof2}, we obtain the following inequality in $\hat{\lambda}_{k+1}U_{k+1}$.
\begin{equation*} 
2\hat{\lambda}_{k+1}U_{k+1} + \frac{(1 - \bar{\rho}^2)}{\bar{\rho}^2}  \hat{\lambda}_{k+1}^2U_{k+1}^2 \leq 1.
\end{equation*}
By solving this inequality, we obtain $\hat{\lambda}_{k+1}U_{k+1} \leq \frac{\bar{\rho}}{1+ \bar{\rho}}$ for all $k \geq 0$.  

Now, by utilizing the resulting inequality and \eqref{eq:B_k}, for all $k \geq 0$, we can show that
\begin{equation*} 
\hat{\lambda}_{k+1} \leq \frac{\bar{\rho}}{(1 + \bar{\rho})U_{k+1}} \leq \frac{2\bar{\rho}}{1 + \bar{\rho}} \cdot \frac{\hat{\lambda}_k}{1 + \hat{\lambda}_k}.  
\end{equation*}
Since $\hat{\lambda}_k = \frac{1}{\hat{\varphi}_k + 1}$, the last inequality leads to 
\begin{equation*}
\hat{\varphi}_{k+1} \geq \frac{1 + \bar{\rho} }{2 \bar{\rho} }\hat{\varphi}_k + \frac{1}{ \bar{\rho}} 
\qquad\Leftrightarrow \qquad 
\hat{\varphi}_{k+1} + \frac{2}{1 - \bar{\rho}} \geq \frac{1 + \bar{\rho}}{2\bar{\rho}}\Big(\hat{\varphi}_k + \frac{2}{1 - \bar{\rho} }\Big).
\end{equation*}
If we denote by $\hat{\psi}_k := \hat{\varphi}_k + \frac{2}{1 - \bar{\rho}}$, then the last inequality means that $\hat{\psi}_{k+1} \geq \big(\frac{1 + \bar{\rho}}{2\bar{\rho}} \big) \hat{\psi}_k$ for all $k\geq 0$.
By induction, we obtain $\hat{\psi}_k \geq \big(\frac{1 + \bar{\rho}}{2\bar{\rho}}\big)^k \hat{\psi}_0$.
Substituting $\hat{\varphi}_k = \hat{\psi}_k - \frac{2}{1 - \bar{\rho} }$ into this bound, we obtain (i).

(ii) 
Since $\bar{\rho} \norms{\hat{\lambda}_ks^k - r^k}  \leq \hat{\lambda}_k\norms{s^k}$ holds for all $k \geq 0$, Lemma~\ref{le:key_properties}(ii) still applies.
Similar to the proof of (i) above, we have $U_{k+1} \geq \frac{1 + \hat{\lambda}_k}{2\hat{\lambda}_k} \geq 1$ from \eqref{eq:B_k} for all $k \geq 0$ due to $0 < \hat{\lambda}_k < 1$.
Moreover,  since $s^0 = r^0$, $U_0 := \frac{\iprods{r^0, s^0}}{\norms{r^0}^2} = 1$.
We conclude that $U_k \geq 1$ for all $k \geq 0$.

The condition $2\omega\hat{\lambda}_k\iprods{s^k, r^k} \leq \norms{r^k}^2$ is equivalent to $\hat{\lambda}_kU_k \leq \frac{1}{2\omega}$.
Since it holds for all $k \geq 1$, we have $U_{k+1} \leq \frac{1}{2\omega\hat{\lambda}_{k+1}}$ for all $k \geq 0$.
Combining this inequality and $U_{k+1} \geq \frac{1 + \hat{\lambda}_k}{2\hat{\lambda}_k}$, we obtain $\frac{1}{\hat{\lambda}_{k+1}} \geq \frac{\omega(1 + \hat{\lambda}_k)}{\hat{\lambda}_k}$.
Substituting $\frac{1}{\hat{\lambda}_k} = \hat{\varphi}_k + 1$ into this inequality, we get
\begin{equation*}
\hat{\varphi}_{k+1} + 1 \geq \omega(\hat{\varphi}_k + 2) \quad \Leftrightarrow \quad \hat{\varphi}_{k+1} + \frac{2\omega-1}{\omega - 1} \geq \omega\left(\hat{\varphi}_k +  \frac{2\omega-1}{\omega - 1} \right).
\end{equation*}
By induction, we have $\hat{\varphi}_k \geq \omega^k\big(\hat{\varphi}_0 + \frac{2\omega-1}{\omega-1}\big) -  \frac{2\omega-1}{\omega - 1}$, which proves (ii).
\end{proof}
%%% End of Proof.

Our next objective is to derive an implementable update rule for $\hat{\lambda}_k$ that does not require the exact contraction factor $\rho$ of $T$, but only an upper bound $\bar{\rho}$ of $\rho$. 
At the same time, the resulting update rule should preserve the linear convergence guarantee of the proposed method. 
The following lemma establishes the key ingredients needed to derive such an update rule.

%%% Lemma 3.
\begin{lemma}\label{le:implementable_step}
Let $\sets{x^k}_{k \geq 0}$ be updated by \eqref{eq:HP_iteration2} and $\sets{\hat{\varphi}_k} \subset (0, +\infty)$ be such that $\hat{\lambda}_k = \frac{1}{1 + \hat{\varphi}_k}$.
Let $r^k$ and $s^k$ be defined by \eqref{eq:residuals}.
Then, the following statements hold.
\begin{compactitem}
\item[$\mathrm{(i)}$]
If we choose $\hat{\lambda}_0 \geq \frac{\bar{\rho}}{1 + \bar{\rho}}$, then $\bar{\rho} \norms{\hat{\lambda}_0s^0 - r^0}  \leq \hat{\lambda}_0\norms{s^0}$ holds.

\item[$\mathrm{(ii)}$]
If $\rho_0 :=  \frac{\norms{x^1 - T(x^0)}}{\norms{x^1 - x^0}}$, then for any $\hat{\lambda}_0 \in \big[\frac{ \bar{\rho}}{1 + \bar{\rho}}, \frac{1}{2} \big)$, we have $\rho_0 = \frac{\hat{\lambda}_0}{1-\hat{\lambda}_0}$ and $\rho \leq \bar{\rho} \leq \rho_0 < 1$.

\item[$\mathrm{(iii)}$] If $r^k \neq 0$ for all $k \geq 0$ and $\hat{\lambda}_{k-1} \in (0, 1)$ for all $k \geq 1$, then $s^k \neq 0$ for all $k \geq 0$.

\item[$\mathrm{(iv)}$] For all $k \geq 1$, let $\hat{\lambda}_k$ and $\rho_k$ respectively be updated by
\begin{equation}\label{eq:lambda_k_update}
\arraycolsep=0.2em
\begin{array}{lcl}
\hat{\lambda}_k & := & \frac{\rho_{k-1}\norms{r^k}^2 }{ \rho_{k-1} \iprods{r^k, s^k} + \sqrt{\rho_{k-1}^2\iprods{r^k, s^k}^2 + (1 - \rho_{k-1}^2 )\norms{r^k}^2\norms{s^k}^2} }, \vspace{1ex}\\
\rho_k  & := & \max\set{\rho_{k-1}, \tau_k}, \quad \textrm{where} \  \tau_k := \frac{\norms{T(x^{k+1}) - T(x^k)}}{\norms{x^{k+1} - x^k}} \ \textrm{or} \ \tau_k := \frac{\norms{x^{k+1} - T(x^k)}}{\norms{x^{k+1} - x^k}},
\end{array}
\end{equation}	
then both $\rho_k \geq \bar{\rho}$ and $\bar{\rho} \norms{\hat{\lambda}_ks^k - r^k}  \leq \hat{\lambda}_k\norms{s^k}$ hold.

\item[$\mathrm{(v)}$]
If we define $\omega := \frac{1 + \rho_0}{2\rho_0}$, then $\omega > 1$ and $2\omega\hat{\lambda}_k\iprods{s^k, r^k} \leq \norms{r^k}^2$ for all $k \geq 1$.
\end{compactitem}
\end{lemma}

%%% Proof of Lemma 3.
\begin{proof}
(i). Since $s^0 = r^0$, the condition $\bar{\rho} \norms{\hat{\lambda}_0s^0 - r^0}  \leq \hat{\lambda}_0\norms{s^0}$ is equivalent to $\bar{\rho}\vert\hat{\lambda}_0 - 1\vert \leq \hat{\lambda}_0$.
However, because $\hat{\lambda}_0 \in (0, 1]$, the last condition is equivalent to $\hat{\lambda}_0 \geq \frac{\bar{\rho}}{1 + \bar{\rho}}$.
 
 (ii). Substituting $x^1 = \hat{\lambda}_0x^0 + (1-\hat{\lambda}_0)T(x^0)$ from \eqref{eq:HP_iteration2} and $s^0 = r^0$ into $\rho_0 =  \frac{\norms{x^1 - T(x^0)}}{\norms{x^1 - x^0}}$ we get $\rho_0 = \frac{\hat{\lambda}_0}{1-\hat{\lambda}_0}$.
 Since $\hat{\lambda}_0 \in \big[\frac{ \bar{\rho}}{1 + \bar{\rho}}, \frac{1}{2}\big)$, we can easily check that that $\rho_0  < 1$.
 Moreover, since $\bar{\rho} \norms{\hat{\lambda}_0s^0 - r^0}  \leq \hat{\lambda}_0\norms{s^0}$ by (i), we conclude that Lemma~\ref{le:key_properties}(i) holds at $k = 0$, i.e., $\rho_0 =  \frac{\norms{x^1 - T(x^0)}}{\norms{x^1 - x^0}} \geq \bar{\rho}$.
 
 (iii)~We prove this statement by induction.
 At $k = 0$, since $r^0 \neq 0$, we have $\norms{s^0} = \norms{r^0} = \norms{x^0 - T(x^0)} > 0$.
 Assume that for any $k \geq 1$, we have $s^{k-1} \neq 0$, we will prove that $s^k \neq 0$.
 Indeed, from the definitions of $s^k$ and $r^k$, and \eqref{eq:HP_iteration2}, we have 
\begin{equation*}
\arraycolsep=0.2em
\begin{array}{lcl}
s^k &= & x^0 - T(x^k) = x^0 - x^k + x^k - T(x^k) = x^0 - \hat{\lambda}_{k-1}x^0 - (1-\hat{\lambda}_{k-1})T(x^{k-1}) + r^k \vspace{1ex}\\
& = & (1-\hat{\lambda}_{k-1})(x^0 - T(x^{k-1})) + r^k =  (1-\hat{\lambda}_{k-1})s^{k-1} + r^k.
\end{array}
\end{equation*}
Thus, we have $\norms{s^k}^2 = (1-\hat{\lambda}_{k-1})^2\norms{s^{k-1}}^2 + 2(1-\hat{\lambda}_{k-1})\iprods{s^{k-1}, r^k} + \norms{r^k}^2$.
Combining this relation and Lemma~\ref{le:key_properties}(ii), we have
\begin{equation*}
\arraycolsep=0.2em
\begin{array}{lcl}
\norms{s^k}^2 & \geq &  (1-\hat{\lambda}_{k-1})^2\norms{s^{k-1}}^2 + \frac{1-\hat{\lambda}_{k-1}}{\hat{\lambda}_{k-1}}\norms{r^k}^2 + \norms{r^k}^2 = (1-\hat{\lambda}_{k-1})^2\norms{s^{k-1}}^2 + \frac{1}{\hat{\lambda}_{k-1}}\norms{r^k}^2.
\end{array}
\end{equation*}
Since $r^{k-1} \neq 0$, $s^{k-1} \neq 0$, and $\hat{\lambda}_{k-1} \in (0, 1)$, the last inequality implies that $s^k \neq 0$.

(iv)~We prove this statement by induction.
 For $k = 1$, from (ii), we have $\rho_0 \geq \bar{\rho}$.
 Using this fact, we conclude that $\bar{\rho} \norms{\hat{\lambda}_1s^1 - r^1}  \leq \hat{\lambda}_1\norms{s^1}$ holds if impose that $\rho_0 \norms{\hat{\lambda}_1s^1 - r^1}  = \hat{\lambda}_1\norms{s^1}$.
The latter equation becomes $(1-\rho_0^2)\hat{\lambda}_1^2\norms{s^1}^2 + 2\rho_0\hat{\lambda}_1\iprods{s^1, r^1} - \rho_0\norms{r^1}^2 = 0$.
Clearly, since $\norms{s^1} > 0$,  $\hat{\lambda}_1$ given by \eqref{eq:lambda_k_update} is the positive solution of this equation. 
Therefore, it satisfies such an equation, and we obtain $\bar{\rho} \norms{\hat{\lambda}_1s^1 - r^1}  \leq \hat{\lambda}_1\norms{s^1}$.
This statement and the fact that $\rho_0 \geq \bar{\rho}$ imply that (iv) holds for $k = 1$.

Suppose that (iv) holds for any $k > 1$, i.e.,  $\rho_{k} \geq \bar{\rho}$ and $\bar{\rho} \norms{\hat{\lambda}_ks^k - r^k}  \leq \hat{\lambda}_k\norms{s^k}$.
We prove that it holds with $k + 1$.

First, by the update rule \eqref{eq:lambda_k_update} and the induction assumption $\rho_k \geq \bar{\rho}$, we have $\rho_{k+1} \geq \rho_k \geq \bar{\rho}$.

Second, since $\rho_k \geq \bar{\rho}$, if we impose $\rho_k \norms{\hat{\lambda}_{k+1}s^{k+1} - r^{k+1}}  = \hat{\lambda}_{k+1}\norms{s^{k+1}}$ then $\bar{\rho} \norms{\hat{\lambda}_{k+1}s^{k+1} - r^{k+1}}  \leq \hat{\lambda}_{k+1}\norms{s^{k+1}}$.
Clearly, as shown in the case $k=1$, $\hat{\lambda}_{k+1}$ computed by \eqref{eq:lambda_k_update} guarantees that $\rho_k \norms{\hat{\lambda}_{k+1}s^{k+1} - r^{k+1}}  = \hat{\lambda}_{k+1}\norms{s^{k+1}}$.
Thus, we conclude that (iv) holds for $k + 1$.

 (v). To guarantee $2\omega\hat{\lambda}_k\iprods{s^k, r^k} \leq \norms{r^k}^2$ for all $k \geq 1$, we need to choose $\omega$ such that
\begin{equation*}
\omega \leq \omega_k := \frac{\norms{r^k}^2}{2\hat{\lambda}_k\iprods{s^k, r^k}} = \frac{1}{2} + \frac{\sqrt{\rho_{k-1}^2\iprods{r^k, s^k}^2 + (1 - \rho_{k-1}^2 )\norms{r^k}^2\norms{s^k}^2} }{2\rho_{k-1}\iprods{s^k, r^k}}.
\end{equation*}
Since $\iprods{s^k, r^k} \leq \norms{s^k}\norms{r^k}$ by the Cauchy-Schwarz inequality, the last expression shows that $\omega_k \geq \frac{1}{2} + \frac{1}{2\rho_{k-1}}$.
Note that the update rule \eqref{eq:lambda_k_update} implies that $\rho_k \leq \rho_0$ for all $k \geq 0$.
Therefore, we have $\omega_k \geq \frac{1}{2} + \frac{1}{2\rho_{k-1}} \geq \frac{1+\rho_0}{2\rho_0}$.
Clearly, since $\rho_0 < 1$, if we choose $\omega :=  \frac{1 + \rho_0}{2\rho_0}$, then $\omega > 1$ and $2\omega\hat{\lambda}_k\iprods{s^k, r^k} \leq \norms{r^k}^2$ holds for all $k \geq 1$.
\end{proof}
%%% End of Proof.

Finally, since $\hat{\lambda}_0 = \frac{1}{\hat{\varphi}_0 + 1}$, if we choose $\hat{\varphi}_0 \in \big(1, \frac{1}{ \bar{\rho} } \big]$, then $\frac{\bar{\rho}}{1 + \bar{\rho}} \leq \hat{\lambda}_0 < \frac{1}{2}$.
Moreover, we also have $\rho_0 = \frac{\hat{\lambda}_0}{1-\hat{\lambda}_0} = \frac{1}{\hat{\varphi}_0}$.
Combining these facts, Lemma~\ref{le:implementable_step}, and the scheme \eqref{eq:HP_iteration2}, we can describe the following implementable version of the adaptive Halpern fixed-point iteration in Algorithm~\ref{alg:adaptive_HP}.

%%% Algorithm 2.
\begin{algorithm}[!htbp]
\caption{(Adaptive Halpern Fixed-Point Algorithm)}\label{alg:adaptive_HP}
\begin{algorithmic}[1]
\State\textbf{Initialization:} Choose an arbitrary initial point $x^0\in\mathcal H$.
\State\label{alg2:step2} Choose an initial value $\hat{\varphi}_0 \in \big(1, \frac{1}{ \bar{\rho} }\big]$.
Compute $\hat{\lambda}_0 : = \frac{1}{\hat{\varphi}_0 + 1}$ and set $\rho_0 := \frac{1}{\hat{\varphi}_0}$.
\State Update $x^1 :=  \hat{\lambda}_0x^0 + (1 - \hat{\lambda}_0) T(x^0)$.
\State \textbf{For $k = 1, 2, \cdots$, perform}
\State\hspace{3ex}Compute $r^k := x^k - T(x^k)$ and $s^k := x^0 - T(x^k)$.
\State\hspace{3ex}If ~$\norms{r^k} = 0$,  then TERMINATE.
\State\hspace{3ex}\label{alg2:step7}Compute $\hat{\lambda}_k \in (0, 1)$ as follows: 
\begin{equation*} 
\hat{\lambda}_k := \frac{\rho_{k-1}\norms{r^k}^2 }{ \rho_{k-1} \iprods{r^k, s^k} + \sqrt{\rho_{k-1}^2\iprods{r^k, s^k}^2 + (1 - \rho_{k-1}^2 )\norms{r^k}^2\norms{s^k}^2} }.
\end{equation*}	
\State\hspace{3ex}\label{alg2:step8}Update
\begin{equation*}
x^{k+1} :=  \hat{\lambda}_kx^0 + (1 - \hat{\lambda}_k) T(x^k).
\end{equation*}
\State\hspace{3ex}\label{alg2:step9}Update $\rho_{k} := \max\left\{ \rho_{k-1}, \tau_k \right\}$, where $\tau_k := \frac{\norms{T(x^{k+1}) - T(x^k)} }{ \norms{x^{k+1} - x^k} }$ or $\tau_k := \frac{\norms{x^{k+1} - T(x^k)} }{ \norms{x^{k+1} - x^k} }$.
\State\textbf{End For}
\end{algorithmic}
\end{algorithm}

\noindent
Let us clarify several aspects of Algorithm~\ref{alg:adaptive_HP}.
\begin{compactitem}
\item Step~\ref{alg2:step2} requires an upper bound $\bar{\rho}$ of the contraction factor $\rho$ in order to define the interval $\big(1, \frac{1}{\bar{\rho}}\big]$ for selecting $\hat{\varphi}_0$. 
If such an upper bound is available, we simply choose $\hat{\varphi}_0 := \frac{1}{\bar{\rho}}$. 
Otherwise, a practical choice is to set $\hat{\varphi}_0 := 1 + \epsilon$ for a sufficiently small $\epsilon > 0$, for example, $\epsilon := 10^{-6}$.

\item Step~\ref{alg2:step7} adds only a negligible computational overhead to the algorithm. 
The main computations consist of one inner product $\iprods{s^k, r^k}$ and two vector norms $\norms{s^k}$ and $\norms{r^k}$, all of which require only $\BigOs{p}$ elementary operations.

\item Step~\ref{alg2:step8} requires the computation of the two norms $\norms{T(x^{k+1}) - T(x^k)}$ and $\norms{x^{k+1} - x^k}$, together with a few basic arithmetic operations. For improved numerical stability, we recommend replacing the update by first computing $\tau_k := \frac{\norms{T(x^{k+1}) - T(x^k)}}{\norms{x^{k+1} - x^k}}$, and then setting $\rho_k := \min\sets{1,\max\sets{\rho_{k-1},\tau_k}}$.
Since $T$ is $\rho$-contractive, we always have $\tau_k \leq \rho$. Consequently, if the initial estimate satisfies $\rho_0 \geq \bar{\rho}$, then $\rho_k=\rho_0$ for all $k\ge0$. 
In practice, however, $\rho_0$ is typically chosen according to \eqref{eq:omega_0_choice}, which may result in $\rho_0\leq\rho$. In this case, the update rule in Step~\ref{alg2:step8} generates a nondecreasing sequence $\rho_0\leq\rho_1\leq\cdots\leq\rho_k\leq\rho$, which provides progressively more accurate estimates of the contraction factor.
\end{compactitem}
Finally, when $\rho=\bar{\rho}=1$, the update of $\hat{\lambda}_k$ in Step~\ref{alg2:step7} reduces to $\hat{\lambda}_k=\frac{\norms{r^k}^2}{2\iprods{s^k,r^k}}$,
which coincides exactly with the parameter update in \cite[Algorithm~3.1]{he2024convergence} for nonexpansive mappings.

%%% 4.2. Convergence analysis.
\subsection{Convergence analysis}\label{subsec:A2_convergence}
We establish linear convergence rates of Algorithm~\ref{alg:adaptive_HP} in the following theorem.
	
%%% Theorem 4.1.	
\begin{theorem}\label{th:convergence_of_adaptive_HP}
Let $ T: \Hil \to \Hil$ be a $\rho$-contractive mapping such that $\mathrm{Fix}(T) \neq \emptyset$ and $\rho \in (0, 1)$.
Let $\bar{\rho}$ be an upper bound of $\rho$ such that $\rho \leq \bar{\rho} < 1$.
Let  $\sets{x^k}_{k\geq 0}$ be generated by Algorithm \ref{alg:adaptive_HP} such that $\hat{\varphi}_0 \in \big(1, \frac{1}{\bar{\rho}}\big]$. 
Then, for all $k \geq 0$, we have 
\begin{equation}\label{eq:convergence_rate}
\norms{ x^k - T(x^k) } \leq \frac{2}{C_k} \norms{x^0 - x^{\star} } \quad \text{and} \quad \norms{x^k - x^{\star}} \leq \frac{2}{(1-\rho)C_k} \norms{x^0 - x^{\star} },
\end{equation}
where $C_k := \big(\frac{1 + \hat{\varphi}_0}{2}\big)^k \big(\hat{\varphi}_0 + \frac{2\hat{\varphi}_0}{\hat{\varphi}_0 -1}\big) - \frac{\hat{\varphi}_0 +  1}{\hat{\varphi}_0 - 1}$ and $x^{\star} \in \mathrm{Fix}(T)$.
Consequently, $\sets{x^k}$ strongly converges to the unique fixed-point $x^{\star}$.
\end{theorem}

%%% Proof of Theorem 4.1.
\begin{proof}
First, by Lemma~\ref{le:implementable_step}(iv), we have $\omega := \frac{1 + \rho_0}{2\rho_0} > 1$.
Moreover, Lemma~\ref{le:implementable_step} shows that both conditions of Lemma~\ref{le:varphi_k_bounds}(ii) are satisfied.
Thus, we conclude that $\hat{\varphi}_k \geq \omega^k\big(\hat{\varphi}_0 + \frac{2\omega-1}{\omega-1}\big) -  \frac{2\omega-1}{\omega - 1}$.
Since $\omega = \frac{1 + \rho_0}{2\rho_0}$ and $\rho_0 = \frac{1}{\hat{\varphi}_0}$, we get $\hat{\varphi}_k + 1 \geq \big( \frac{1 + \hat{\varphi}_0}{2} \big)^k\big(\hat{\varphi}_0 + \frac{2\hat{\varphi}_0}{\hat{\varphi}_0-1}\big) - \frac{\hat{\varphi}_0 + 1}{\hat{\varphi}_0 - 1}$.
Moreover, since $\hat{\varphi}_0 \in \big(1, \frac{1}{ \bar{\rho} }\big]$, it is obvious to check that $\hat{\lambda}_0 = \frac{1}{1 + \hat{\varphi}_0} \in \big[\frac{\bar{\rho}}{1 + \bar{\rho}}, \frac{1}{2}\big)$.

Next, following the proof in \cite[Theorem 3.2]{he2024convergence}, it is easy to verify that the equality:
\begin{equation*} 
\varphi \norms{a}^2 + 2\iprods{a, b} + \norms{ c }^2 -\norms{a + c}^2 = \tfrac{\varphi+1}{2}\norms{ a}^2 - \tfrac{2}{\varphi+1}\norms{a + c - b}^2 + \frac{2}{\varphi+1}\norms{ a + c - b - \tfrac{\varphi+1}{2}a}^2
\end{equation*}
holds for all $a, b, c \in \mathcal{H}$ and an arbitrary positive number $\varphi$. 

Substituting $\varphi := \hat{\varphi}_k$, $a:= x^k - T(x^k)$, $b:= x^k - x^0$, and $c:=T(x^k) - x^{\star}$, and the facts that $a + c = x^k-x^{\star}$, $a + c - b=x^0-x^{\star}$ into the last equality, we obtain
\begin{equation}\label{eq:th41_proof1}
\arraycolsep=0.2em
\begin{array}{lcl}
\Tc_{[1]} & := & \hat{\varphi}_k \norms{ x^k - T(x^k) }^2 + 2\iprods{x^k - Tx^k, x^k - x^0 } + \norms{T(x^k) - x^{\star}}^2 - \norms{x^k - x^{\star} }^2 \vspace{1ex} \\
& = & \frac{\hat{\varphi}_k+1}{2}\norms{ x^k - T(x^k) }^2 - \frac{2}{\hat{\varphi}_k+1}\norms{x^0 - x^{\star} }^2 + \frac{2}{ \hat{\varphi}_k+1}\norms{x^0 - x^{\star} - \frac{\hat{\varphi}_k+1}{2}(x^k - T(x^k) )}^2.
\end{array}
\end{equation}
Combining \eqref{eq:th41_proof1} and  Lemma \ref{le:key_properties} (iii), we get
\begin{equation*}
\frac{\hat{\varphi}_k+1}{2}\norms{x^k - T(x^k)}^2 - \frac{2}{\hat{\varphi}_k+1}\norms{ x^0 - x^{\star} }^2 \leq 0.
\end{equation*}
which implies that
\begin{equation*}
\norms{ x^k - T(x^k)}  \leq \frac{2}{\hat{\varphi}_k+1}\norms{ x^0-x^{\star}}.
\end{equation*}
Substituting $\hat{\varphi}_k + 1 \geq C_k := \big( \frac{1 + \hat{\varphi}_0}{2} \big)^k\big(\hat{\varphi}_0 + \frac{2\hat{\varphi}_0}{\hat{\varphi}_0-1}\big) - \frac{\hat{\varphi}_0 + 1}{\hat{\varphi}_0 - 1}$ into the last expression, we obtain the first bound in \eqref{eq:convergence_rate}, where $\frac{1}{C_k} \to 0$ geometrically as $k \to \infty$.

Finally, the proof of the second bound in \eqref{eq:convergence_rate} is similar to the proof in Theorem~\ref{th:convergence_theorem1}, and we omit it.
Here, since $\hat{\varphi}_0 > 1$, we notice that $\frac{1}{C_k}$ converges linearly to zero as $k \to \infty$.
\end{proof}
%%% End of Proof.


%%%%%%%%%%%%%%%%%%%%%%%%%%%%%%%%%%%%%%%%%%%%%%%%%%%%%%%%%%%%%%%%
%%% 5. Applications to Nonexpansive Mappings via Tikhonov's Regularization
%%%%%%%%%%%%%%%%%%%%%%%%%%%%%%%%%%%%%%%%%%%%%%%%%%%%%%%%%%%%%%%%
\section{Applications to Co-Coercive Equations}\label{sec:NonExMapping}
In this section, we apply the results developed in Sections~\ref{sec:linear_convergence_rate} and \ref{sec:adaptive_Halpern_method} to solve the co-coercive equation \eqref{eq:CE}. The main idea is summarized as follows.
\begin{compactitem}
\item First, we apply Tikhonov regularization to the operator $G$ in \eqref{eq:CE} by defining $G_{\mu} := G + \mu \Id$ for a sufficiently small parameter $\mu > 0$, where $\Id$ denotes the identity operator.

\item Next, we appropriately choose a parameter $\eta > 0$ and define $T_{\mu} := \Id - \eta G_{\mu}$ so that $T_{\mu}$ becomes a contractive mapping.

\item We then apply either Algorithm~\ref{alg:adaptive_HP0} from Section~\ref{sec:linear_convergence_rate} or Algorithm~\ref{alg:adaptive_HP} from Section~\ref{sec:adaptive_Halpern_method} to approximate the fixed point $x^{\star}_{\mu}$ of $T_{\mu}$.

\item Finally, we characterize an approximate solution of \eqref{eq:CE} by choosing an appropriate value of the regularization parameter $\mu$.
\end{compactitem}
Notice that the fixed-point problem \eqref{eq:FP} with a nonexpansive mapping $T$ is a special case of \eqref{eq:CE}. 
Indeed, by defining $G := \Id - T$, solving \eqref{eq:FP} is equivalent to solving \eqref{eq:CE}, where $G$ is $\frac{1}{2}$-co-coercive. 
Therefore, it suffices to consider the more general problem \eqref{eq:CE} throughout this section.

%%% 5.1. Tikhonov's regularization for \eqref{eq:CE}
\subsection{Tikhonov regularization for \eqref{eq:CE}}\label{subsec:Tikhonov_reg}
Let $G : \Hil \to \Hil$ be a single-valued mapping and $\beta$-co-coercive, i.e., $\iprods{G(x) - G(y), x - y} \geq \beta\norms{G(x) - G(y)}^2$ for all $x, y \in \Hil$, and $\mu > 0$ be a given positive number.
We consider two families of regularized mappings $G_{\mu}$ of $G$ and $T_{\mu}$ respectively defined as follows:
\begin{equation}\label{eq:T_mu}
G_{\mu} := G + \mu \Id 
\quad \text{and} \quad T_{\mu} := \Id - \eta G_{\mu} = (1-\eta\mu)\Id - \eta G,
\end{equation}
for an appropriate $\eta > 0$ determined later.
First, we prove the following lemma.

%% Lemma 5.1.
\begin{lemma}\label{le:regularized_mapping}
Let $G : \Hil \to \Hil$ be a $\beta$-co-coercive mapping and $\zer{G} \neq \emptyset$.
Let $G_{\mu}$ and $T_{\mu}$ be defined by \eqref{eq:T_mu} for some $\mu > 0$ and $\eta > 0$.
Then, the following properties hold:
\begin{compactitem}
\item[$\mathrm{(i)}$] $x^{\star}_{\mu} \in \zer{G_{\mu}}$ iff $x^{\star}_{\mu} \in \mathrm{Fix}(T_{\mu})$ for any $\eta > 0$ and $\mu > 0$.
\item[$\mathrm{(ii)}$] For $x^{\star} \in \zer{G}$ and $x^{\star}_{\mu} \in \zer{G_{\mu}}$, we have 
\begin{equation}\label{eq:G_mu_pro1}
\mu^2\norms{x^{\star}_{\mu} - x^{\star}}^2 + (1 + 2\mu\beta)\norms{G(x^{\star}_{\mu}) - G(x^{\star})}^2 \leq \mu^2\norms{x^{\star}}^2.
\end{equation}
Consequently, we have $\norms{G(x_{\mu}^{\star})} \leq \frac{\mu}{\sqrt{1+2\mu\beta}}\norms{x^{\star}}$ and $\norms{x^{\star}_{\mu} - x^{\star}} \leq \norms{x^{\star}}$.
\item[$\mathrm{(iii)}$] If we choose $\eta$ such that $0 < \eta \leq \frac{2\beta}{1 + 2\mu\beta}$, then $T_{\mu}$ is $\rho_{\mu}$-contractive with $\rho_{\mu} := 1 - \mu\eta \in (0, 1)$.
\end{compactitem}
\end{lemma}

%%% Proof of Lemma 5.1.
\begin{proof}
\noindent{$\mathrm{(i)}$}
We note that $x^{\star}_{\mu}\in \zer{G_{\mu}}$ means that $G_{\mu}(x^{\star}_{\mu}) = 0$, which is equivalent to $x^{\star}_{\mu} = x^{\star}_{\mu} - \eta G_{\mu}(x^{\star}_{\mu}) = T_{\mu}(x^{\star}_{\mu})$ for any $\mu > 0$ and $\eta > 0$. This means that $x^{\star}_{\mu}$ is a fixed-point of $T_{\mu}$.

\noindent{$\mathrm{(ii)}$}
Let $x^{\star} \in \zer{G}$ and $x^{\star}_{\mu} \in \zer{G_{\mu}}$.
Then, we have $-\mu x^{\star} = G(x^{\star}_{\mu} ) - G( x^{\star} ) + \mu(x^{\star}_{\mu} - x^{\star})$.
Squaring both sides, we obtain 
\begin{equation*}
\mu^2\norms{x^{\star}}^2 = \norms{G(x^{\star}_{\mu} ) - G( x^{\star} )}^2 + 2\mu\iprods{G(x^{\star}_{\mu} ) - G( x^{\star} ), x^{\star}_{\mu} - x^{\star}} + \mu^2\norms{x^{\star}_{\mu} - x^{\star}}^2.
\end{equation*}
By the $\beta$-co-coercivity of $G$, this relation leads to \eqref{eq:G_mu_pro1}.
The two remaining bounds are direct consequences of \eqref{eq:G_mu_pro1} and $G(x^{\star}) = 0$.

\noindent{$\mathrm{(iii)}$}~Since $0 < \eta \leq \frac{2\beta}{1 + 2\beta\mu}$, we have $\eta - 2\beta(1 - \eta\mu) \leq 0$.
In addition, since $\frac{2\beta}{1 + 2\beta\mu} < \frac{1}{\mu}$, we get $1 - \eta\mu > 0$.
Now, using these facts, $G_{\mu}$ and $T_{\mu}$ defined by \eqref{eq:T_mu}, and the $\beta$-co-coercivity of $G$, we can show that
\begin{equation*}
\arraycolsep=0.2em
\begin{array}{lcl}
    \norms{ T_{\mu} (x) - T_{\mu} (y) }^2 &= &  \norms{ (1-\eta \mu) (x-y) - \eta (G(x) - G(y)) }^2 \vspace{1ex}\\
    &= & (1-\eta \mu )^2 \norms{ x - y }^2 + \eta^2 \norms{ G(x) - G(y) }^2 - 2\eta (1-\eta \mu) \iprods{ G(x) - G(y), x - y } \vspace{1ex} \\    
    &\leq & (1 - \eta \mu )^2 \norms{ x - y }^2 + \eta^2 \norms{ G(x) - G(y) }^2 - 2\eta (1-\eta \mu) \beta \norms{ G(x) - G(y) }^2 \vspace{1ex} \\
    & = & (1 - \eta \mu)^2 \norms{ x - y }^2 + \eta [ \eta - 2\beta (1-\eta \mu ) ] \norms{ G(x) - G(y) }^2 \vspace{1ex}\\
    & \leq & (1 - \eta\mu)^2\norms{x - y}^2.
\end{array}
\end{equation*}
Since $1 - \eta\mu > 0$, the last expression leads to $ \norms{ T_{\mu} (x) - T_{\mu} (y) } \leq (1 - \eta\mu)\norms{x - y}$, showing that $T_{\mu}$ is $\rho_{\mu}$-contractive with $\rho_{\mu} := 1 - \eta\mu$.
\end{proof}
%%% End of Proof.

Lemma~\ref{le:regularized_mapping} provides useful insights for developing practical algorithms.
First, the bound $\norms{G(x_{\mu}^{\star})} \leq \frac{\mu}{\sqrt{1+2\mu\beta}}\norms{x^{\star}}$ in Lemma~\ref{le:regularized_mapping}(ii) shows that $\norms{G(x^{\star}_{\mu})}$ converges to zero with the same rate as $\mu$ when $\mu \to 0^{+}$.
Moreover, since $G$ is continuous and $\sets{x^{\star}_{\mu}}$ is bounded due to the estimate $\norms{x^{\star}_{\mu} - x^{\star}} \leq \norms{x^{\star}}$, we conclude that $x^{\star} = \lim_{\mu\to 0^{+}}x^{\star}_{\mu}$ (or a cluster of $\sets{x^k}$) is a solution of \eqref{eq:CE}.
Second, if the co-coercivity parameter $\beta$ is known, then we may choose $\eta := \frac{2\beta}{1 + 2\beta\mu}$.
With this choice, the mapping $T_{\mu}$ depends only on the regularization parameter $\mu$.
In practice, however, $\beta$ is often unavailable. In this case, we may estimate it by $\beta_0 : = \frac{\iprods{G(x^0 + v) - G(x^0), v}}{\norms{G(x^0 + v) - G(x^0)}^2}$ for some nonzero direction $v \in\Hil$.

%%% 5.2. Parameter-free Halpern methods for co-coercive equations
\subsection{Parameter-free Halpern methods for co-coercive equations}\label{subsec:methods_for_CE}
For simplicity, we present the application of Algorithm~\ref{alg:adaptive_HP0} to approximate a fixed point of $T_{\mu}$. 
The extension to Algorithm~\ref{alg:adaptive_HP} is very analogous. 
Once an approximate fixed point of $T_{\mu}$ is obtained, we recover an approximate solution of \eqref{eq:CE} by appropriately choosing the regularization parameter $\mu$.

Since the co-coercivity parameter $\beta$ of $G$ is generally unknown, we estimate it by the upper bound
\begin{equation}\label{eq:beta_0_estimate}
0 < \beta \leq \beta_0 := \frac{\iprods{G(x^0 + v) - G(x^0), v}}{\norms{G(x^0 + v) - G(x^0)}^2},
\end{equation}
for a given direction $0 \neq v\in\Hil$ such that $G(x^0 + v) \neq G(x^0)$.

First, by Lemma~\ref{le:regularized_mapping}(ii), if we choose $\mu := \epsilon > 0$ for a sufficiently small tolerance $\epsilon > 0$, then $x^{\star}_{\mu} = x^{\star}_{\epsilon}$ and $\norms{G(x^{\star}_{\epsilon})} \leq \frac{\norms{x^{\star}}\epsilon}{\sqrt{1 + 2\beta\epsilon}}$, showing that $x^{\star}_{\epsilon}$ is an $\hat{\epsilon}$-approximate solution of \eqref{eq:CE}, where $\hat{\epsilon} := \frac{\epsilon}{\sqrt{1 + 2\beta\epsilon}}$.

Next, after computing $\beta_0$, Lemma~\ref{le:regularized_mapping}(iii) suggests choosing $\eta := \frac{\kappa \beta_0}{1 + 2\beta_0\epsilon}$ for some $\kappa \in (0, 1)$ (e.g., $\kappa := 0.9$).
The purpose of this choice is to ensure that $0 < \eta \leq \frac{2\beta}{1 + 2\mu\beta}$.

With these choices of $\mu$ and $\eta$, we specialize Algorithm~\ref{alg:adaptive_HP0} to approximate a fixed point of $T_{\mu}$, or equivalently, an approximate solution of \eqref{eq:CE}, resulting in Algorithm~\ref{alg:AHP_mu}.

%%%% Algorithm 1.
\begin{algorithm}[!htbp]
\caption{(Parameter-Free HP Algorithm for Co-coercive Equation \eqref{eq:CE})}\label{alg:AHP_mu}
\begin{algorithmic}[1]
\State\textbf{Initialization:} Choose an initial point $x^0\in\mathcal H$ and an accuracy $\epsilon \in (0, 1)$.
\State Choose $\mu := \epsilon$ and $\eta := \frac{\kappa\beta_0}{1 + 2\beta_0\epsilon}$ for $\beta_0$ in \eqref{eq:beta_0_estimate} and $\kappa \in (0, 1)$ (e.g., $\kappa := 0.9$).
\State Estimate $\omega_0$ as in \eqref{eq:omega_0_choice} and set $\varphi_{-1} := 0$.
\State Evaluate $T_{\mu}(x^0) := (1 - \eta\mu)x^0 - \eta G(x^0)$.
\State \textbf{For $k = 0, 1, 2, \cdots$, perform}
\State\hspace{3ex}If ~$\norms{x^k - T_{\mu}(x^k)} = 0$,  then TERMINATE.
\State\hspace{3ex}Update $\bar{\varphi}_{k} := \bar{\varphi}_{k-1} + \omega_k^{2k}$.
\State\hspace{3ex}\label{alg3:step5}Update
\begin{equation*}
x^{k+1} :=  \frac{1}{ \bar{\varphi}_k + 1} x^0 + \frac{ \bar{\varphi}_k}{\bar{\varphi}_k + 1}T_{\mu}(x^k).
\end{equation*}
\State\hspace{3ex}\label{alg3:step6}Evaluate $T_{\mu}(x^{k+1}) := (1-\eta\mu)x^{k+1} - \eta G(x^{k+1})$.
\State\hspace{3ex}\label{alg3:step7}Update $\omega_{k+1}$ as in \eqref{eq:update_omega_k} or \eqref{eq:update_omega_k_modified} with $T_{\mu}$.
\State\textbf{End For}
\end{algorithmic}
\end{algorithm}

Finally, we prove the following corollary, which states the convergence of Algorithm~\ref{alg:AHP_mu}.

%%% Corollary 1.
\begin{corollary}\label{co:convergence_of_AHP}
Let $G : \Hil \to \Hil$ in \eqref{eq:CE} be $\beta$-co-coercive and $\mathrm{zer}(G)\neq\emptyset$.
Suppose that $\mu := \epsilon > 0$, $0 < \eta \leq \frac{2\beta}{1 + 2\beta\epsilon}$, and $x^{\star} \in \mathrm{zer}(G)\neq\emptyset$.
Let $\sets{x^k_{\mu}}_{k\geq 0}$ be generated by Algorithm~\ref{alg:AHP_mu}.
Then, we have 
\begin{equation}\label{eq:AHP_mu_bound} 
\arraycolsep=0.2em
\begin{array}{lcl}
\norms{ G(x^k_{\mu}) }  & \leq & \frac{6(1 - \eta\epsilon)^{k-1}}{\eta^2\epsilon^2}\big(\norms{G(x^0)} + \epsilon\norms{x^0}\big)  +  2\epsilon\norms{x^{\star}}.
\end{array}
\end{equation}
Consequently, the maximum number of iterations $k$ to achieve  $\norms{ G(x^k_{\mu}) } \leq M_0\epsilon$ is $k := \BigOs{\epsilon^{-1}\ln(\epsilon^{-1})}$, where $M_0 := \frac{6}{\eta^2}\big(\norms{G(x^0)} + \norms{x^0}\big) + 2\norms{x^{\star}}$.

\end{corollary}

%% Proof of Corollary 1.
\begin{proof}
Since $x^k_{\mu}$ is generated by Algorithm~\ref{alg:AHP_mu}, by \eqref{eq:convergence_bound1} and $\eta G_{\mu} = \Id - T_{\mu}$, we have
\begin{equation*} 
\eta\norms{ G_{\mu}(x^k_{\mu}) } = \norms{x^k_{\mu} - T_{\mu}(x_{\mu}^k)} \leq  \Big[ 1 + \frac{1+\rho_{\mu}}{(1-\rho_{\mu})^2\rho_{\mu}}\Big] \rho_{\mu}^k \norms{x^0 - T_{\mu}(x^0)} \leq \frac{3\eta \rho_{\mu}^{k-1} }{(1-\rho_{\mu})^2}\norms{G_{\mu}(x^0)},
\end{equation*}
where we have used the fact that $1 + \frac{1+\rho_{\mu}}{(1-\rho_{\mu})^2\rho_{\mu}} \leq \frac{3}{(1-\rho_{\mu})^2\rho_{\mu}}$.

Since $\rho_{\mu} := 1 - \eta\mu = 1 - \eta\epsilon$, this expression leads to 
\begin{equation*} 
\norms{ G_{\mu}(x^k_{\mu}) } \leq \frac{3(1 - \eta\epsilon)^{k-1}}{\eta^2\epsilon^2}\norms{G_{\mu}(x^0)}.
\end{equation*}
Moreover, following the same proof for the second bound of \eqref{eq:convergence_bound1}, we have  $\norms{x^k_{\mu} - x^{\star}_{\mu}} \leq \frac{\eta}{1 - \rho_{\mu}}\norms{G_{\mu}(x^k)}$.
Substituting $\rho_{\mu} = 1 - \eta\mu = 1 - \eta\epsilon$, this inequality leads to
\begin{equation*} 
\norms{ x^k_{\mu} -  x^{\star}_{\mu} } \leq \frac{3(1 - \eta\epsilon)^{k-1}}{\eta^2\epsilon^3}\norms{G_{\mu}(x^0)}.
\end{equation*}
This inequality together with $\norms{x_{\mu}^{\star} - x^{\star}}  \leq \norms{x^{\star}}$ from Lemma~\ref{le:regularized_mapping}(ii) show that
\begin{equation*} 
\arraycolsep=0.2em
\begin{array}{lcl}
\norms{x^k_{\mu}} & \leq & \norms{x^k_{\mu} - x_{\mu}^{\star}} + \norms{x_{\mu}^{\star} - x^{\star}} + \norms{x^{\star}}  \leq \frac{3(1 - \eta\epsilon)^{k-1}}{\eta^2\epsilon^3}\norms{G_{\mu}(x^0)} + 2\norms{x^{\star}}.
\end{array}
\end{equation*}
Combining the last two inequalities, we can show that
\begin{equation*} 
\arraycolsep=0.2em
\begin{array}{lcl}
\norms{ G(x^k_{\mu}) } & \leq & \frac{3(1 - \eta\epsilon)^{k-1}}{\eta^2\epsilon^2}\norms{G_{\mu}(x^0)} + \epsilon\norms{x_{\mu}^k} \leq \frac{3(1+\eta)(1 - \eta\epsilon)^{k-1}}{\eta^2\epsilon^2}\norms{G_{\mu}(x^0)}  + 2 \epsilon \norms{x^{\star}} \vspace{1ex}\\
& \leq & \frac{6(1 - \eta\epsilon)^{k-1}}{\eta^2\epsilon^2 }\big(\norms{G(x^0)} + \epsilon\norms{x^0}\big)  +  2\epsilon\norms{x^{\star}}.
\end{array}
\end{equation*}
For a sufficiently small $\epsilon \in (0, 1)$, we have $\eta := \frac{\kappa \beta_0}{1 + 2\beta_0\epsilon} > \frac{\kappa\beta_0}{2}$.
This means that $\eta = \BigOs{1}$.

Since we have chosen $\mu = \epsilon > 0$, by choosing $k \geq 1 + \frac{3\ln(\epsilon)}{\ln(1 - \eta\epsilon)} = \BigOs{\epsilon^{-1}\ln(\epsilon^{-1})}$, we can show that $\frac{(1 - \eta\epsilon)^{k-1}}{\epsilon^2} \leq \epsilon$.
Therefore, the last expression leads to
\begin{equation*} 
\arraycolsep=0.2em
\begin{array}{lcl}
\norms{ G(x^k_{\mu}) } & \leq & \Big[\frac{6}{\eta^2}\big(\norms{G(x^0)} + \norms{x^0}\big) + 2\norms{x^{\star}}\Big] \epsilon = M_0\epsilon,
\end{array}
\end{equation*}
where $M_0 := \frac{6}{\eta^2}\big(\norms{G(x^0)} + \norms{x^0}\big) + 2\norms{x^{\star}}$.
This proves \eqref{eq:AHP_mu_bound}.
Consequently, we conclude that the worst-case iteration complexity of Algorithm~\ref{alg:AHP_mu} is $\BigOs{\epsilon^{-1}\ln(\epsilon^{-1})}$.
\end{proof}
%%% End of Proof.

Note that Algorithm~\ref{alg:AHP_mu} has an iteration complexity of $\BigOs{\epsilon^{-1}\ln(\epsilon^{-1})}$.
This bound is only slightly worse than the $\BigOs{\epsilon^{-1}}$ complexity achieved by the standard Halpern fixed-point iteration with the theoretical parameter choice $\varphi_k = k$, as stated in Theorem~\ref{th:linear_convergence_of_HP}, by a logarithmic factor of $\epsilon^{-1}$.

%%% Remark 5.1.
\begin{remark}\label{re:generalization}
Several important classes of problems can be reformulated as \eqref{eq:CE}, or equivalently as the fixed-point problem \eqref{eq:FP}, including the optimality conditions of convex optimization problems, convex-concave saddle-point problems, and monotone variational inequality problems (VIPs). 
Depending on the underlying problem structure, different reformulations of \eqref{eq:CE} can be employed, such as those based on proximal (or resolvent) operators, forward-backward splitting, backward-forward splitting, Douglas-Rachford splitting, and three-operator splitting. 
Under appropriate assumptions, the resulting reformulation mapping $G$ is co-coercive. Consequently, the algorithms developed in this paper can be directly applied to solve these problems. 
Since these reformulations are well established in the literature, we omit the details and refer the interested reader to the corresponding references.
\end{remark}

%%%%%%%%%%%%%%%%%%%%%%%%%%%%%%%%%%%%%%%%%%%%%%%%%%%%%%%%%%%%%%%%
%%% 6. From Halpern's Iteration to Parameter-Free Nesterov's Acceleration.
%%%%%%%%%%%%%%%%%%%%%%%%%%%%%%%%%%%%%%%%%%%%%%%%%%%%%%%%%%%%%%%%
\section{Parameter-Free Nesterov's Accelerated Fixed-Point Methods}\label{sec:NesMethods}
As we mentioned earlier, since $T$ is $\rho$-contractive, both the Banach-Picard (BP) and Krasnosel'ski\v{\i}-Mann (KM) fixed-point iterations converge linearly to the unique fixed point of $T$. 
However, in practice, it is often difficult to determine whether the underlying mapping is $\rho$-contractive for some $\rho \in (0,1)$ or merely nonexpansive. 
Consequently, it is generally safer to apply the KM fixed-point iteration, which remains convergent under the weaker assumption of nonexpansiveness.

On the other hand, the KM fixed-point iteration achieves only a $\BigOs{1/\sqrt{k}}$ convergence rate for the fixed-point residual $\norms{x^k - T(x^k)}$ in the nonexpansive setting, which is significantly slower than the optimal $\BigOs{1/k}$ rate attained by Halpern's fixed-point iteration. 
Recently, \cite{attouch2020convergence}, and subsequently \cite{bot2022bfast}, proposed Nesterov-type accelerated variants of the KM fixed-point iteration. 
Interestingly, \cite{tran2022connection} showed that these accelerated KM schemes are mathematically equivalent to Halpern's fixed-point iteration \eqref{eq:HP_iteration}, differing only in the choice of the parameter sequence $\lambda_k$. 
Motivated by this observation, we exploit this equivalent representation to develop parameter-free Nesterov-accelerated KM fixed-point methods for approximating a fixed point of $T$.

%%% 6.1. The derivation of the algorithm.
\subsection{The derivation of the algorithm}\label{subsec:Nes_derivation}
To derive Nesterov's accelerated fixed-point schemes for solving \eqref{eq:FP}, we exploit the ideas from \cite{tran2022connection} and process as follows.
First, from \eqref{eq:HP_iteration2}, we can write it at the $(k-1)^{th}$ and the $k^{th}$ iterations:
\begin{equation*}
\arraycolsep=0.2em
\begin{array}{lcl}
(\varphi_{k-1} + 1)x^k & = &  x^0 + \varphi_{k-1}T(x^{k-1}), \vspace{1ex}\\
(\varphi_k + 1)x^{k+1} & = &  x^0 + \varphi_kT(x^k).
\end{array}
\end{equation*}
Next, eliminating $x^0$ in both lines, we can show that $(\varphi_k + 1)x^{k+1} - (\varphi_{k-1} + 1)x^k = \varphi_kT(x^k) - \varphi_{k-1}T(x^{k-1})$.
This equation leads to the following update rule:
\begin{equation}\label{eq:NesFP_scheme0}
\arraycolsep=0.2em
\begin{array}{lcl}
x^{k+1} &= & \frac{\varphi_{k-1} + 1}{\varphi_k + 1}x^k +  \frac{\varphi_k}{\varphi_k + 1}T(x^k) - \frac{\varphi_{k-1}}{\varphi_k + 1}T(x^{k-1}).
\end{array}
\end{equation}
This scheme requires two consecutive iterates $x^{k-1}$ and $x^k$ to update $x^{k+1}$.
Therefore, at the first iteration $k=0$, we update $x^1 := \frac{1}{\varphi_0 + 1}x^0 +  \frac{\varphi_0}{\varphi_0 + 1}T(x^0)$.

Note that we can rewrite \eqref{eq:NesFP_scheme0} as follows:
\begin{equation*} 
\arraycolsep=0.2em
\begin{array}{lcl}
x^{k+1} & := & (1 - \lambda_k)x^k +  \lambda_k T(x^k) + \gamma_k(T(x^k) - T(x^{k-1})).
\end{array}
\end{equation*}
where $\lambda_k := \frac{\varphi_k - \varphi_{k-1}}{\varphi_k + 1}$ and $\gamma_k := \frac{\varphi_{k-1}}{\varphi_k + 1}$.
If $\varphi_k > \varphi_{k-1} \geq 0$, then we have $\lambda_k \in (0, 1]$.
Clearly, if $\gamma_k = 0$, then \eqref{eq:NesFP_scheme0} reduces to the standard KM fixed-point iteration \cite{Bauschke2011}.
Therefore, we can view \eqref{eq:NesFP_scheme0} as a KM fixed-point variant with a \textit{correction term} $\gamma_k(T(x^k) - T(x^{k-1}))$.

Now, we exploit the technique in Algorithm~\ref{alg:adaptive_HP} to construct an update rule for $\varphi_k$ for all $k \geq 0$.
We note that $s^k := x^0 - T(x^k) = (\varphi_{k-1} + 1)x^k - \varphi_{k-1}T(x^{k-1}) - T(x^k) = \varphi_{k-1}(x^k - T(x^{k-1})) + r^k$.
If we denote by $q^k := x^k - T(x^{k-1})$, then we have $s^k = \varphi_{k-1}q^k + r^k$.
The condition $\rho_{k-1}\norms{\hat{\lambda}_ks^k - r^k} = \hat{\lambda}_k\norms{s^k}$ with $\hat{\lambda}_k = \frac{1}{\varphi_k + 1}$ from Lemma~\ref{le:implementable_step} becomes
\begin{equation*} 
\rho_{k-1}^2\norms{\varphi_{k-1}q^k - \varphi_kr^k}^2 = \norms{\varphi_{k-1}q^k + r^k}^2.
\end{equation*}
Solving this quadratic equation, we obtain 
\begin{equation}\label{eq:varphi_k_update}
\varphi_k :=   \frac{\rho_{k-1}\varphi_{k-1}\iprods{q^k, r^k} + \sqrt{\Delta_k}}{\rho_{k-1}\norms{r^k}^2},
\end{equation}
where $\Delta_k := \norms{r^k}^2\norms{\varphi_{k-1}q^k + r^k}^2 - \rho_{k-1}^2\varphi_{k-1}^2\big(\norms{q^k}^2\norms{r^k}^2  - \iprods{q^k, r^k}^2 \big) > 0$.

Finally, similar to Algorithm~\ref{alg:adaptive_HP}, we can update $\rho_k$ as $\rho_k := \max\sets{\rho_{k-1}, \frac{\norms{T(x^{k+1}) - T(x^k)}}{\norms{\Delta{x}^k}}}$, where $\Delta{x}^k := x^{k+1} - x^k$.
Alternatively, we can also compute $\rho_{k} := \max\set{\rho_{k-1}, \frac{\norms{q^{k+1}} }{\norms{\Delta{x}^k} }}$ for $\rho_k$. 

%%%% 6.2. 
\subsection{The algorithm and its convergence}\label{subsec:Nes_convergence}
Now, we are ready to present the complete description of the proposed method in Algorithm~\ref{alg:adaptive_NesFP}.

%%% Algorithm 4.
\begin{algorithm}[!htbp]
\caption{(Parameter-Free Nesterov's Accelerated Fixed-Point Algorithm)}\label{alg:adaptive_NesFP}
\begin{algorithmic}[1]
\State\textbf{Initialization:} Choose an arbitrary initial point $x^0\in\mathcal H$.
\State \textbf{Option 1:} Choose $\omega_0$ as in \eqref{eq:new_stepsize0} and $\varphi_0 := \omega_0$.
\textbf{Option 2:} Choose $\varphi_0 = \frac{1}{\rho_0} \in \big(1, \frac{1}{\rho}\big]$.
\State Update $x^1 := \frac{x^0 + \varphi_0T(x^0)}{\varphi_0 + 1}$.
\State \textbf{For $k = 1, 2, \cdots$, perform}
\State\hspace{3ex}Compute $r^k := x^k - T(x^k)$ and $q^k := x^k - T(x^{k-1})$.
\State\hspace{3ex}If ~$\norms{r^k} = 0$,  then TERMINATE.
\State\hspace{3ex}\textbf{Option 1:} Compute $\varphi_{k} := \varphi_{k-1} + \omega_k^{2k}$.
\State\hspace{3ex}\textbf{Option 2:} Compute $\Delta_k := \norms{r^k}^2\norms{\varphi_{k-1}q^k + r^k}^2 - \rho_{k-1}^2\varphi_{k-1}^2\big(\norms{q^k}^2\norms{r^k}^2  - \iprods{q^k, r^k}^2 \big)$ and 
\begin{equation*}
\varphi_k :=   \frac{\rho_{k-1}\varphi_{k-1}\iprods{q^k, r^k} + \sqrt{\Delta_k}}{\rho_{k-1}\norms{r^k}^2}.
\end{equation*}
\State\hspace{3ex}Update
\begin{equation*}
x^{k+1} =  \frac{\varphi_{k-1} + 1}{\varphi_k + 1}x^k +  \frac{\varphi_k}{\varphi_k + 1}T(x^k) - \frac{\varphi_{k-1}}{\varphi_k + 1}T(x^{k-1}).
\end{equation*}
\State\hspace{3ex}\textbf{Option 1:} Update $\omega_{k+1}$ as in Algorithm~\ref{alg:adaptive_HP0}.
\State\hspace{3ex}\textbf{Option 2:} Update $\rho_k$ as in Algorithm~\ref{alg:adaptive_HP}.
\State\textbf{End For}
\end{algorithmic}
\end{algorithm}

\noindent
We make the following observations.
\begin{compactitem}
\item Each iteration of Algorithm~\ref{alg:adaptive_NesFP} requires only one evaluation of $T(x^k)$. The remaining computations consist of vector additions, scalar-vector multiplications, norm evaluations, and inner products, all of which require $\BigOs{p}$ elementary operations. Hence, the per-iteration computational cost of Algorithm~\ref{alg:adaptive_NesFP} is essentially the same as that of Algorithms~\ref{alg:adaptive_HP0} and \ref{alg:adaptive_HP}.

	\item In contrast to the Halpern-type schemes, including Algorithms~\ref{alg:adaptive_HP0} and \ref{alg:adaptive_HP}, the update $x^{k+1}$ in Algorithm~\ref{alg:adaptive_NesFP} does not explicitly involve the anchor point $x^0$. This feature may help reduce the accumulation of errors associated with repeatedly incorporating the anchor point. Moreover, for nonexpansive mappings, Algorithm~\ref{alg:adaptive_NesFP} may potentially achieve the faster asymptotic rate $o(1/k)$, compared with the $\BigOs{1/k}$ rates typically obtained by Halpern-type fixed-point methods \cite{tran2022connection}.
\end{compactitem}
Finally, since the scheme \eqref{eq:NesFP_scheme0} is mathematically equivalent to \eqref{eq:HP_iteration2}, the convergence of Algorithm~\ref{alg:adaptive_NesFP} follows directly from the convergence results established for Algorithms~\ref{alg:adaptive_HP0} and \ref{alg:adaptive_HP}. The following corollary formalizes this observation.

%%% Corollary 5.1.
\begin{corollary}\label{co:Nes_convergence}
Let $T : \Hil \to \Hil$ be a $\rho$-contractive mapping with $\rho \in (0,1)$ and assume that $\mathrm{Fix}(T)\neq\emptyset$. Let $\{x^k\}_{k\ge0}$ be the sequence generated by Algorithm~\ref{alg:adaptive_NesFP}. Then, all the conclusions of Theorems~\ref{th:convergence_theorem1} and \ref{th:convergence_of_adaptive_HP} continue to hold.
\end{corollary}

%%%%%%%%%%%%%%%%%%%%%%%%%%%%%%%%%%%%%%%%%%%%%%%%%%%%%%%%%
%%% 7. Numerical Examples
%%%%%%%%%%%%%%%%%%%%%%%%%%%%%%%%%%%%%%%%%%%%%%%%%%%%%%%%%
\section{Numerical Experiments}\label{sec:num_examples}
We present a series of numerical experiments to evaluate the performance of the proposed algorithms and compare them with recent existing methods from the literature. 
Specifically, we implement the following nine algorithms.
%%
\begin{compactitem}
\item The geometric Halpern scheme \eqref{eq:HP_iteration} from \cite{park2022exact} with $\varphi_k := \sum_{i=1}^k\rho^{-2i}$, denoted by \texttt{[Geo.~HP]}.

\item Algorithm~3.1 from \cite{he2024convergence}, denoted by \texttt{[Alg.~3.1~(He~et~al.)]}.

\item Two variants of our Algorithm~\ref{alg:adaptive_HP0}, corresponding to the update rules \eqref{eq:update_omega_k} and \eqref{eq:update_omega_k_modified}, denoted by \texttt{[Alg.~1~$(\omega_{\max})$]} and \texttt{[Alg.~1~$(\omega_{\min})$]}, respectively.

\item Two variants of Algorithm~\ref{alg:adaptive_HP}, using the two update rules for $\rho_k$ specified in the algorithm, denoted by \texttt{[Alg.~2~(adpt.)]} and \texttt{[Alg.~2~(adpt2)]}, respectively.

\item Three variants of Algorithm~\ref{alg:adaptive_NesFP}. The geometric variant corresponding to \textbf{Option~1} is denoted by \texttt{[Alg.~3~(geo.)]}. The remaining two variants, corresponding to \textbf{Option~2} with the two update rules for $\rho_k$ used in Algorithm~\ref{alg:adaptive_HP}, are denoted by \texttt{[Alg.~3~(adpt.)]} and \texttt{[Alg.~3~(adpt2)]}, respectively.
\end{compactitem}
%%
All algorithms were implemented in Python (Jupyter Notebook) with assistance from OpenAI Codex.
The implementations were carefully reviewed and independently verified by the authors. All experiments were conducted on a MacBook Pro Apple M4, 24GB memory, and 1Tb storage.
%Our code can be found at \href{https://github.com/unc-optimization}{https://github.com/unc-optimization}.

Each algorithm is terminated once the relative fixed-point residual satisfies $\norms{x^k-T(x^k)} \leq \texttt{tol}\cdot\max\left\{1,\norms{x^0-T(x^0)}\right\}$, where $\texttt{tol}$ is a prescribed tolerance. Unless otherwise specified, the maximum number of iterations is set to $k_{\max} := 10^4$.

%%% 7.1. Experiments on contractive mappings.
\subsection{Experiments on fixed-point problems with contractive mappings}\label{subsec:example1}
We conduct two groups of numerical experiments on fixed-point problems associated with contractive mappings. The first group (\textbf{Example 1.1}) consists of linear contractive mappings, whereas the second group (\textbf{Example 1.2}) considers nonlinear contractive mappings.

%%% Example 1.1.
\vspace{2ex}
\noindent\textbf{Example 1.1.} 
We first consider a linear mapping $T(x) := \mathbf{Q}x + \mathbf{q}$, where $\mathbf{Q} \in \R^{10\times 10}$ and $\mathbf{q}\in\R^{10}$ satisfy $\rho := \norms{\mathbf{Q}}_2 < 1$. 
In our experiments, we randomly generate $n=50$ different mappings $T_i(x)$ with contraction factors $\rho \in [0.5, 0.99]$. 
The vector $\mathbf{q} \in \mathcal{N}(0, 0.5)$ is generated independently from the normal distribution. 
The true solution $x^{\star}$ can be computed explicitly as $x^{\star} := (\Id - \mathbf{Q})^{-1}\mathbf{q}$. 
For all test instances, we use the same initial point $x^0 := \mathrm{ones}(p,1)$.

Table~\ref{tbl:performance_comparison_linear_case_50} summarizes the performance of the $9$ algorithms in reaching an approximate solution $x^k$ with the stopping tolerance $\texttt{tol} = 10^{-8}$. 
Here, \textbf{Median Iter.} and \textbf{Mean Iter.} denote the median and mean numbers of iterations, respectively; \textbf{Median Res.} and \textbf{Mean Res.} report the median and mean residual norms, respectively; and \textbf{Median Dist.} and \textbf{Mean Dist.} represent the median and mean values of $\norms{x^k - x^{\star}}$, respectively.

\begin{table}[!htbp]
\centering
\caption{Performance comparison of $9$ algorithms on $50$ linear contractive mappings}
\label{tbl:performance_comparison_linear_case_50}
\setlength{\tabcolsep}{5pt}
\renewcommand{\arraystretch}{1.12}
\resizebox{\textwidth}{!}{%
\begin{tabular}{l c c c c c c c}
\toprule
\textbf{Method} & \textbf{Successes} & \textbf{Median Iter.} & \textbf{Mean Iter.} & \textbf{Median Res.} & \textbf{Mean Res.} & \textbf{Median Dist.} & \textbf{Mean Dist.} \\
\midrule
Geo. HP & 50 & 25.5 & 74.5 & 3.05e-08 & 3.05e-08 & 4.05e-08 & 4.23e-08 \\
Alg. 3.1 (He et al.) & 0 & 10000.0 & 10000.0 & 8.63e-05 & 9.14e-05 & 1.05e-04 & 1.04e-04 \\
Alg. 1 ($\omega_{\max}$) & 50 & 20.5 & 23.1 & 2.86e-08 & 2.76e-08 & 2.55e-08 & 3.47e-08 \\
Alg. 1 ($\omega_{\min}$) & 50 & 21.0 & 24.1 & 2.47e-08 & 2.65e-08 & 2.54e-08 & 3.55e-08 \\
Alg. 2 (adpt.) & 50 & 28.0 & 28.5 & 2.84e-08 & 2.91e-08 & 3.95e-08 & 4.46e-08 \\
Alg. 2 (adpt2) & 50 & 24.5 & 27.0 & 2.73e-08 & 2.94e-08 & 3.74e-08 & 4.47e-08 \\
Alg. 3 (geo.) & 50 & 21.0 & 23.7 & 2.75e-08 & 2.75e-08 & 2.47e-08 & 3.40e-08 \\
Alg. 3 (adpt.) & 50 & 28.0 & 28.5 & 2.84e-08 & 2.91e-08 & 3.95e-08 & 4.46e-08 \\
Alg. 3 (adpt2) & 50 & 24.5 & 27.0 & 2.73e-08 & 2.94e-08 & 3.74e-08 & 4.47e-08 \\
\bottomrule
\end{tabular}%
}
\end{table}

As shown in Table~\ref{tbl:performance_comparison_linear_case_50}, all variants of the proposed methods perform well on this test problem and consistently outperform \texttt{Geo.~HP} in terms of the number of iterations. 
Among them, \texttt{Alg.~1~$(\omega_{\max})$} and \texttt{Alg.~3~(geo.)} exhibit the best overall performance. 
In contrast, \texttt{Alg.~3.1~(He et al.)} from \cite{he2024convergence} fails to converge within the maximum limit of $10^4$ iterations, reaching only an approximate solution with an accuracy of approximately $10^{-4}$. 
This behavior is expected, as the method in \cite{he2024convergence} is designed for general nonexpansive mappings, whereas the methods proposed in this paper are specifically tailored to exploit the stronger contractivity assumption.

Figure~\ref{fig:lin_maps_50} further illustrates the performance of the seven successful algorithms. 
The left panel shows the distribution of the numbers of iterations required for convergence, while the right panel reports the final residual attained by each algorithm on each of the $50$ test mappings.

\begin{figure}[!h]
	\centering
	\includegraphics[width=\textwidth]{figs/example_1_distributions}
	\caption{Performance comparison of the $8$ algorithms to solve \eqref{eq:FP} on $50$ linear mappings}
	\label{fig:lin_maps_50}
\end{figure}

%%% Example 1.2.
\vspace{2ex}
\noindent\textbf{Example 1.2.} 
Next, we perform a similar experiment to \textbf{Example~1.1}, but with nonlinear mappings of the form $T(x) :=  \alpha\tanh(\mathbf{Q}x + \mathbf{q})) + \mathbf{f}$.
To construct a contractive mapping $T$, we proceed as follows. We first generate $\mathbf{Q} \in \R^{10\times 10}$ with entries drawn independently from the standard normal distribution and then rescale it so that $\norms{\mathbf{Q}}_2 \in [0.5, 1.5]$. 
Next, we choose $\alpha = \frac{\rho}{\norms{\mathbf{Q}}_2}$ for $\rho \in [0.5, 0.99]$. 
Finally, the vectors $\mathbf{q}$ and $\mathbf{f}$ are generated independently from the standard normal distribution. 
By construction, the resulting mapping $T$ is $\rho$-contractive with $\rho \in [0.5, 0.99]$.

Using this procedure, we generate $50$ random mappings and evaluate the same $9$ algorithms considered in \textbf{Example~1.1}. 
All experimental settings are identical to those in \textbf{Example~1.1}. 
The results are summarized in Table~\ref{tbl:performance_comparison_nonlinear_case_50}, which reports the same performance metrics as Table~\ref{tbl:performance_comparison_linear_case_50}.

\begin{table}[!htbp]
\centering
\caption{Performance comparison of $9$ algorithms on $50$ nonlinear contractive mappings}
\label{tbl:performance_comparison_nonlinear_case_50}
\setlength{\tabcolsep}{5pt}
\renewcommand{\arraystretch}{1.12}
\resizebox{\textwidth}{!}{%
\begin{tabular}{l c c c c c c c}
\toprule
\textbf{Method} & \textbf{Successes} & \textbf{Median Iter.} & \textbf{Mean Iter.} & \textbf{Median Res.} & \textbf{Mean Res.} & \textbf{Median Dist.} & \textbf{Mean Dist.} \\
\midrule
Geo. HP & 50 & 36.0 & 60.8 & 3.39e-08 & 3.62e-08 & 3.48e-08 & 3.67e-08 \\
Alg. 3.1 (He et al.) & 0 & 10000.0 & 10000.0 & 8.63e-05 & 8.97e-05 & 9.51e-05 & 9.37e-05 \\
Alg. 1 ($\omega_{\max}$) & 50 & 14.0 & 14.7 & 2.79e-08 & 2.92e-08 & 2.00e-08 & 2.10e-08 \\
Alg. 1 ($\omega_{\min}$) & 50 & 14.0 & 14.8 & 2.40e-08 & 2.56e-08 & 1.69e-08 & 2.03e-08 \\
Alg. 2 (adpt.) & 50 & 16.0 & 16.9 & 2.87e-08 & 2.94e-08 & 3.13e-08 & 2.96e-08 \\
Alg. 2 (adpt2) & 50 & 15.0 & 15.9 & 2.32e-08 & 2.60e-08 & 2.27e-08 & 2.51e-08 \\
Alg. 3 (geo.) & 50 & 13.0 & 14.2 & 2.31e-08 & 2.79e-08 & 1.74e-08 & 2.10e-08 \\
Alg. 3 (adpt.) & 50 & 16.0 & 16.9 & 2.87e-08 & 2.94e-08 & 3.13e-08 & 2.96e-08 \\
Alg. 3 (adpt2) & 50 & 15.0 & 15.9 & 2.32e-08 & 2.60e-08 & 2.27e-08 & 2.51e-08 \\
\bottomrule
\end{tabular}%
}
\end{table}

Again, we observe performance trends similar to those in \textbf{Example~1.1}. 
All proposed methods successfully solve the nonlinear test problems, with \texttt{Alg.~1~$(\omega_{\max})$} and \texttt{Alg.~3~(geo.)} remaining among the best-performing methods in terms of the number of iterations. 
Figure~\ref{fig:nonlin_maps_50} further illustrates the performance of the eight successful algorithms. 
The left panel shows the distribution of the numbers of iterations required for convergence, while the right panel reports the final residual attained by each algorithm on each of the $50$ nonlinear mappings.

\begin{figure}[!h]
	\centering
	\includegraphics[width=\textwidth]{figs/example_2_distributions}
	\caption{Performance comparison of the $7$ successful algorithms to solve \eqref{eq:FP} on $50$ nonlinear mappings}
	\label{fig:nonlin_maps_50}
\end{figure}

%%%% 7.2. Experiments on fixed-point problems of nonexpansive mapping.
\subsection{Experiments on fixed-point problems of nonexpansive mapping}\label{subsec:exam2}
In this section, we evaluate the performance of the proposed algorithms on fixed-point problems involving nonexpansive mappings. 
Specifically, we compare the nine algorithms developed in this paper with the adaptive Halpern method of \cite{he2024convergence} and the classical Halpern fixed-point iteration using the step size $\lambda_k = \frac{1}{k+1}$. 
Since this iteration is equivalent to the Halpern method in \cite{park2022exact} when the contraction factor is unavailable, we continue to denote it by \texttt{Geo.~HP}.

%%% Example 2.1.
\vspace{2ex}
\noindent\textbf{Example 2.1.}
We consider a nonexpansive mapping $T : \R^3 \to \R^3$ studied in \cite{he2021optimal,he2024convergence}, which is explicitly defined as follows:
\begin{equation}\label{eq:T_exam3}
T(x) := \frac{1}{54.5}\begin{pmatrix}
-35x_1 - \sqrt{|x_1|+1} - 10x_2 + 14x_3 + 1\\
-10x_1 - 26x_2 - 0.5\sin(x_2) + 4x_3\\
14x_1 + 4x_2 - 38x_3 - \arctan(0.5x_3).
\end{pmatrix}
\end{equation}
As claimed in  \cite{he2021optimal,he2024convergence}, this mapping is nonexpansive. 
Moreover, its fixed-point is $x^{\star} = (0, 0, 0)^{\top}$.

Now, we apply the $9$ algorithms described above to approximate a fixed point of this mapping. 
We choose the initial point $x^0 = (1, 1, 1)^{\top}$ and run each algorithm for at most $50,\!000$ iterations. 
The algorithms are terminated once the relative residual reaches the prescribed tolerance $\mathrm{tol} = 10^{-8}$. 
Figure~\ref{fig:Exam51} illustrates the performance of the algorithms, where the left panel shows the evolution of the fixed-point residual $\norms{x^k - T(x^k)}$ and the right panel shows the distance to the solution $\norms{x^k - x^{\star}}$ as functions of the iteration count.

\begin{figure}[!h]
	\centering
	\includegraphics[width=\textwidth]{figs/paper_example_51_convergence}
	\caption{The convergence progress of the $9$ algorithms to approximate a fixed point of $T$ in \eqref{eq:T_exam3}}
	\label{fig:Exam51}
\end{figure}

As shown in Figure~\ref{fig:Exam51}, \texttt{Geo.~HP}, \texttt{Alg.~3.1.~(He et al.)}, \texttt{Alg.~2~(adpt.)}, and \texttt{Alg.~3~(adpt.)} exhibit similar convergence behavior during the early stages of the iterations. 
However, when the algorithms are run for a sufficiently large number of iterations, \texttt{Alg.~3.1.~(He et al.)} becomes the slowest among these methods. 
The method \texttt{Alg.~1~$(\omega_{\min})$} performs noticeably better, while the remaining four proposed methods achieve substantially faster convergence.
Among all the methods tested, \texttt{Alg.~1~$(\omega_{\max})$} and \texttt{Alg.~3~(geo.)} deliver the best overall performance.

%%% Example 2.2.
\vspace{2ex}
\noindent\textbf{Example 2.2.} 
Following \cite{he2024convergence}, we evaluate the proposed algorithms on the well-known LASSO (Least Absolute Shrinkage and Selection Operator) problem. 
In particular, we consider the same experimental setup as Example~5.2 in \cite{he2024convergence}, where the problem is formulated as follows:
\begin{equation}\label{eq:lasso}
\min_{x \in \R^p} \Big\{\Lc(x) := \tfrac{1}{2m}\norms{Ax - b}^2_2 + \tau \norms{x}_1 \Big\},
\end{equation}
where $A \in \R^{m\times p}$ and $b \in \R^m$ are the input data, and $\tau > 0$ is a regularization parameter.

If we define $T(x) : = \mathrm{prox}_{\gamma\tau\norms{\cdot}_1}\big(x - \gamma A^{\top}(Ax - b) \big)$, then $T$ is nonexpansive, where $\gamma > 0$ and $\mathrm{prox}_f$ is the proximal operator of $f$.
Solving \eqref{eq:lasso} is equivalent to finding a fixed-point $x^{\star}$ of $T$.

To test our algorithms and their competitors, we use $10$ LASSO problems with dimensions $p = 512j$ for $j = 1,\ldots,10$, and $m = \lceil0.6p\rceil$. 
For each problem, we generate $30$ random instances.
The design matrix $A$ is generated from a correlated Gaussian model with AR(1) covariance (i.e., autoregressive of order $1$), where the correlation coefficient is set to $0.3$, and each column is normalized to have the $\ell_2$-norm $\sqrt{m}$. 
The ground-truth solution $x^{\sharp}$ is chosen to be sparse with approximately $5\%$ nonzero entries, whose magnitudes are sampled from a shifted Gaussian distribution to avoid arbitrarily small coefficients. 
The observation vector is generated according to $y=Ax^\star+\varepsilon$, where the noise vector $\varepsilon$ is scaled to achieve a signal-to-noise ratio of $25$ dB. 
For each problem instance, the regularization parameter is selected as $\tau = 0.05\,\tau_{\max}$ with $\tau_{\max}=\frac{\norms{A^\top y}_\infty}{m}$, where $\tau_{\max}$ is the smallest value of $\tau$ for which the zero vector is an optimal solution to \eqref{eq:lasso}.

We evaluate the $8$ adaptive algorithms described above on $300$ LASSO instances spanning $10$ different problem dimensions. 
For all experiments, we use the initial point $x^0 = \mathrm{ones}(p,1)$. 
Each algorithm is terminated once either the maximum number of $5000$ iterations is reached or the relative error satisfies the prescribed tolerance $\mathrm{tol} = 10^{-6}$.

The aggregate performance of the $8$ algorithms is reported in Table~\ref{tbl:aggregate_performance_lasso}. 
Here, \textbf{Obj.} and \textbf{Time (s)} denote the mean objective value and the mean CPU time (in seconds), respectively. 
The remaining performance metrics are the same as those reported in Table~\ref{tbl:performance_comparison_linear_case_50}.

\begin{table}[!htbp]
\centering
\small
\caption{Overall performance of $8$ algorithms on $30$ instances of $10$ different LASSO problems.}
\label{tbl:aggregate_performance_lasso}
\setlength{\tabcolsep}{5pt}
\renewcommand{\arraystretch}{1.12}
\resizebox{\textwidth}{!}{%
\begin{tabular}{lrrrrrrrr}
\toprule
\textbf{Method} {\!\!\!\!\!} & {\!\!\!\!} \textbf{Success} {\!\!\!} & {\!\!\!} \textbf{Median Iter.} {\!\!\!} & {\!\!\!} \textbf{Mean Iter.} {\!\!\!} & {\!\!\!} \textbf{Median Rel. Res.} {\!\!\!} & {\!\!\!} \textbf{Mean Res.} {\!\!\!} & {\!\!\!} \textbf{Max Res.} {\!\!\!} & {\!\!\!} \textbf{Obj.} {\!\!\!} & {\!\!\!} \textbf{Time[s]} {\!\!\!} \\
\midrule
Alg. 3.1 [He et al.] & 0/300 & 5000.0 & 5000.0 & 1.94e-06 & 8.63e-05 & 8.63e-05 & 52.34899 & 5.5314 \\
Alg. 1 (omega max) & 300/300 & 78.0 & 78.9 & 9.20e-07 & 3.94e-05 & 5.98e-05 & 52.34816 & 0.0904 \\
Alg. 1 (omega min) & 300/300 & 84.0 & 87.7 & 9.29e-07 & 3.96e-05 & 6.00e-05 & 52.34825 & 0.0957 \\
Alg. 2 (adpt.) & 300/300 & 96.0 & 97.3 & 9.32e-07 & 3.99e-05 & 5.99e-05 & 52.34845 & 0.1102 \\
Alg. 2 (adpt2) & 300/300 & 92.0 & 92.3 & 9.26e-07 & 3.96e-05 & 5.96e-05 & 52.34841 & 0.1053 \\
Alg. 3 (geo.) & 300/300 & 78.0 & 79.0 & 9.25e-07 & 3.95e-05 & 6.10e-05 & 52.34816 & 0.0901 \\
Alg. 3 (adpt.) & 300/300 & 96.0 & 97.3 & 9.32e-07 & 3.99e-05 & 5.99e-05 & 52.34845 & 0.1105 \\
Alg. 3 (adpt2) & 300/300 & 92.0 & 92.3 & 9.26e-07 & 3.96e-05 & 5.96e-05 & 52.34841 & 0.1056 \\
\bottomrule
\end{tabular}%
}
\end{table}


\begin{figure}[!h]
	\centering
	\includegraphics[width=\textwidth]{figs/paper_example_52_distributions}
	\caption{The distribution of iteration numbers $($left$)$ and residuals $($right$)$ of 8 algorithms on LASSO problem}
	\label{fig:aggregate_performance_lasso}
\end{figure}

As shown in Table~\ref{tbl:aggregate_performance_lasso}, \texttt{Alg.~3.1~(He~et~al.)} fails to achieve the prescribed relative residual tolerance of $10^{-6}$ within the maximum limit of $5000$ iterations. 
In contrast, all the proposed algorithms successfully attain the target accuracy in fewer than $100$ iterations on average. 
Consistent with the previous experiments, \texttt{Alg.~1~$(\omega_{\max})$} and \texttt{Alg.~3~(geo.)} exhibit the best overall performance.

Figure~\ref{fig:aggregate_performance_lasso} further illustrates the performance of the algorithms. 
The left panel shows the distribution of the numbers of iterations required for convergence, while the right panel reports the final residual attained by each algorithm over all problem instances. 
We observe that the numbers of iterations remain remarkably stable across different problem instances and problem dimensions, indicating that the proposed methods are robust with respect to the problem size.
 
 %%% Example 2.3.
\vspace{2ex}
\noindent\textbf{Example 2.3.} 
We perform experiments on a collection of $50$ nonexpansive mappings defined by
\begin{equation*}
T(x) := \frac{1}{2}\big( \mathrm{proj}_C(x) + Q\tanh(x) \big), \quad \textrm{where}\quad C := [-1, 1]^p \ \textrm{and} \ Q^{\top}Q = \Id_p.
\end{equation*}
Here, $\mathrm{proj}_C$ denotes the projection onto the closed and convex set $C$. 
It is straightforward to verify that $T$ is nonexpansive but not globally contractive. 
Moreover, $T$ has the unique fixed point $x^{\star} = 0$.

To generate the test instances, we randomly construct $50$ mappings of this form in $\R^{10}$ by sampling random orthonormal matrices $Q$. 
We then apply the $9$ algorithms to these $50$ instances. 
Each algorithm is run for at most $30,\!000$ iterations and is terminated once the prescribed tolerance $\mathrm{tol} = 10^{-5}$ is achieved. 
For all experiments, we use the initial point $x^0 = \mathrm{ones}(p,1)$.

The aggregate performance of the algorithms is summarized in Table~\ref{tbl:performance_comparison_nonexpansive_mapping_50}, where the reported performance metrics are the same as those in Table~\ref{tbl:performance_comparison_linear_case_50}. 
Figure~\ref{fig:nonlin_nonexpansive_maps_50} provides a complementary visualization of the results by showing the distribution of the numbers of iterations (left) and the final relative residuals (right) for the algorithms that successfully converged.

\begin{table}[!htbp]
\centering
\caption{Performance comparison of $9$ algorithms on $50$ nonlinear nonexpansive mappings}
\label{tbl:performance_comparison_nonexpansive_mapping_50}
\setlength{\tabcolsep}{5pt}
\renewcommand{\arraystretch}{1.12}
\resizebox{\textwidth}{!}{%
\begin{tabular}{l c c c c c c c}
\toprule
\textbf{Method} & \textbf{Successes} & \textbf{Median Iter.} & \textbf{Mean Iter.} & \textbf{Median Res.} & \textbf{Mean Res.} & \textbf{Median Dist.} & \textbf{Mean Dist.} \\
\midrule
Geo. HP & 0 & 30000.0 & 30000.0 & 1.05e-04 & 1.05e-04 & 8.29e-02 & 8.22e-02 \\
Alg. 3.1 (He et al.) & 0 & 30000.0 & 30000.0 & 8.63e-05 & 8.63e-05 & 7.69e-02 & 7.63e-02 \\
Alg. 1 ($\omega_{\max}$) & 50 & 1633.0 & 1526.8 & 2.70e-05 & 2.50e-05 & 8.74e-02 & 8.44e-02 \\
Alg. 1 ($\omega_{\min}$) & 50 & 2944.5 & 2617.1 & 2.70e-05 & 2.55e-05 & 7.74e-02 & 7.57e-02 \\
Alg. 2 (adpt.) & 50 & 2280.0 & 2271.1 & 2.70e-05 & 2.56e-05 & 8.08e-02 & 7.81e-02 \\
Alg. 2 (adpt2) & 50 & 1703.5 & 1625.5 & 2.70e-05 & 2.56e-05 & 8.69e-02 & 8.40e-02 \\
Alg. 3 (geo.) & 50 & 1633.5 & 1528.9 & 2.69e-05 & 2.57e-05 & 8.74e-02 & 8.44e-02 \\
Alg. 3 (adpt.) & 50 & 2280.0 & 2271.1 & 2.70e-05 & 2.56e-05 & 8.08e-02 & 7.81e-02 \\
Alg. 3 (adpt2) & 50 & 1703.5 & 1625.5 & 2.70e-05 & 2.56e-05 & 8.69e-02 & 8.40e-02 \\
\bottomrule
\end{tabular}%
}
\end{table}

\begin{figure}[!h]
	\centering
	\includegraphics[width=\textwidth]{figs/example_3_distributions}
	\caption{Performance comparison of the $7$ algorithms to solve \eqref{eq:FP} on $50$ nonlinear nonexpansive mappings}
	\label{fig:nonlin_nonexpansive_maps_50}
\end{figure}

As shown in Table~\ref{tbl:performance_comparison_nonexpansive_mapping_50}, \texttt{Geo.~HP}, which uses the step size $\lambda_k = \frac{1}{k+1}$, does not achieve the prescribed tolerance $\mathrm{tol} = 10^{-6}$ within the maximum limit of $30,\!000$ iterations. 
Similarly, \texttt{Alg.~3.1.~(He~et~al.)}, whose theoretical convergence rate is $\BigOs{1/k}$, also fails to attain the target accuracy within the allotted iterations. 
In contrast, all the proposed methods successfully reach the prescribed tolerance, requiring fewer than $3000$ iterations on average. 
Among them, \texttt{Alg.~1~$(\omega_{\max})$} and \texttt{Alg.~3~(geo.)} exhibit the best overall performance, while the remaining proposed methods perform comparably.

This experiment further demonstrates that the proposed algorithms remain effective for nonexpansive mappings, even though they are primarily designed for contractive mappings. 
As expected, their convergence is generally slower than in the contractive setting, especially when applying to this challenging example, reflecting the weaker assumptions on the underlying mapping.

%%%%% Acknowledgements.
\vspace{1ex}
\noindent\textbf{Acknowledgements.}
This work was partially supported by the National Science Foundation (NSF) through the RTG grant DMS-2134107 and by the Office of Naval Research (ONR) under grant No.~N00014-23-1-2588. 
A substantial portion of this work was completed while the first two authors, Quoc Tran-Dinh and Pham Ngoc Anh, were visiting the Vietnam Institute for Advanced Study in Mathematics (VIASM), Hanoi, Vietnam, in June 2026.

%%%%% Declaration of using AI tools
\vspace{1ex}
\noindent\textbf{Declaration of AI-Assisted Tools.}
During the preparation of this manuscript, we used large language models, including ChatGPT and Gemini, to assist in verifying elementary mathematical derivations, algebraic manipulations, and properties of mathematical expressions. We also used Codex to assist with implementing the proposed algorithms, generating synthetic data, and preparing numerical experiments.
%
All mathematical proofs, algorithmic implementations, numerical results, and conclusions presented in this paper were carefully reviewed and independently verified by the authors. The authors take full responsibility for the correctness and integrity of the content of this manuscript.

%%%%%%%%%%%%%%%%%%%%%%%%%%%%%%%%%%%%%%%%%%%%%%%%%%%%%%%%
%%%% Bibliographies 
%%%%%%%%%%%%%%%%%%%%%%%%%%%%%%%%%%%%%%%%%%%%%%%%%%%%%%%%
\bibliographystyle{plain}
%\bibliography{/Users/quoc/MyDropboxFolder/E-Books/tran_bibtex_2025}
%\end{document}

\begin{thebibliography}{10}
\itemsep=0.0em
\bibitem{alakoya2021modified}
T.~Alakoya, L.~Jolaoso, and O.~T. Mewomo.
\newblock Modified inertial subgradient extragradient method with self adaptive
  stepsize for solving monotone variational inequality and fixed point
  problems.
\newblock {\em Optimization}, 70(3):545--574, 2021.

\bibitem{attouch2020convergence}
H.~Attouch and A.~Cabot.
\newblock Convergence of a relaxed inertial proximal algorithm for maximally
  monotone operators.
\newblock {\em Math. Program.}, 184(1):243--287, 2020.

\bibitem{Banach1922}
S.~Banach.
\newblock Sur les opérations dans les ensembles abstraits et leur application
  aux équations intégrales.
\newblock {\em Fundamenta Mathematicae}, 3:133--181, 1922.

\bibitem{Barzilai1988}
J.~Barzilai and J.~M. Borwein.
\newblock {Two-Point Step Size Gradient Methods}.
\newblock {\em IMA J Numer Anal}, 8(1):141--148, January 1988.

\bibitem{Bauschke2011}
H.~H. Bauschke and P.~Combettes.
\newblock {\em Convex analysis and monotone operators theory in {H}ilbert
  spaces}.
\newblock Springer-Verlag, 2nd edition, 2017.

\bibitem{bot2022bfast}
R.~I. Bot and D.~K. Nguyen.
\newblock Fast {K}rasnosel\'skii-{M}ann algorithm with a convergence rate of
  the fixed point iteration of $o(1/k)$.
\newblock {\em SIAM J. Numer. Anal.}, 61(6):2813--2843, 2023.

\bibitem{ioan2023relaxed}
R.~I. Bot, M.~Sedlmayer, and P.~T. Vuong.
\newblock A relaxed inertial forward-backward-forward algorithm for solving
  monotone inclusions with application to {GANs}.
\newblock {\em J. Mach. Learn. Res. (JMLR)}, 24:1--37, 2023.

\bibitem{carmon2022making}
Y.~Carmon and O.~Hinder.
\newblock Making {SGD} parameter-free.
\newblock In {\em Conference on learning theory}, pages 2360--2389. PMLR, 2022.

\bibitem{colao2015krasnoselskii}
V.~Colao and G.~Marino.
\newblock {K}rasnoselskii-{M}ann method for non-self mappings.
\newblock {\em Fixed Point Theory and Applications}, 2015(1):39, 2015.

\bibitem{daskalakis2018training}
C.~Daskalakis, A.~Ilyas, V.~Syrgkanis, and H.~Zeng.
\newblock Training {GANs} with {O}ptimism.
\newblock In {\em International Conference on Learning Representations (ICLR
  2018)}, pages 1--9, 2018.

\bibitem{defazio2022parameter}
A.~Defazio and K.~Mishchenko.
\newblock Parameter free dual averaging: {O}ptimizing {L}ipschitz functions in
  a single pass.
\newblock In {\em OPT 2022: Optimization for Machine Learning (NeurIPS 2022
  Workshop)}, 2022.

\bibitem{defazio2023learning}
A.~Defazio and K.~Mishchenko.
\newblock Learning-rate-free learning by d-adaptation.
\newblock In {\em International conference on machine learning}, pages
  7449--7479. PMLR, 2023.

\bibitem{diakonikolas2020halpern}
J.~Diakonikolas.
\newblock {H}alpern iteration for near-optimal and parameter-free monotone
  inclusion and strong solutions to variational inequalities.
\newblock In {\em Conference on Learning Theory}, pages 1428--1451. PMLR, 2020.

\bibitem{Duchi2011}
J.~Duchi, E.~Hazan, and Y.~Singer.
\newblock Adaptive subgradient methods for online learning and stochastic
  optimization.
\newblock {\em J. Mach. Learn. Res.}, 12:2121--2159, 2011.

\bibitem{Facchinei2003}
F.~Facchinei and J.-S. Pang.
\newblock {\em Finite-dimensional variational inequalities and complementarity
  problems}, volume 1-2.
\newblock Springer-Verlag, 2003.

\bibitem{fercoq2019adaptive}
O.~Fercoq and Z.~Qu.
\newblock Adaptive restart of accelerated gradient methods under local
  quadratic growth condition.
\newblock {\em IMA Journal of Numerical Analysis}, 39(4):2069--2095, 2019.

\bibitem{halpern1967fixed}
B.~Halpern.
\newblock Fixed points of nonexpanding maps.
\newblock {\em Bull. Am. Math. Soc.}, 73(6):957--961, 1967.

\bibitem{he2021optimal}
S.~He, Q.-L. Dong, H.~Tian, and X.-H. Li.
\newblock On the optimal relaxation parameters of {K}rasnosel'ski--{M}ann
  iteration.
\newblock {\em Optimization}, 70(9):1959--1986, 2021.

\bibitem{he2018totally}
S.~He, T.~Wu, A.~Gibali, and Q.-L. Dong.
\newblock Totally relaxed, self-adaptive algorithm for solving variational
  inequalities over the intersection of sub-level sets.
\newblock {\em Optimization}, 67(9):1487--1504, 2018.

\bibitem{he2024convergence}
S.~He, H.-K. Xu, Q-L. Dong, and N.~Mei.
\newblock Convergence analysis of the {H}alpern iteration with adaptive
  anchoring parameters.
\newblock {\em Math. Comput.}, 93(345):327--345, 2024.

\bibitem{ito2023parameter}
M.~Ito, Z.~Lu, and C.~He.
\newblock A parameter-free conditional gradient method for composite
  minimization under {H}{\"o}lder condition.
\newblock {\em Journal of Machine Learning Research}, 24(166):1--34, 2023.

\bibitem{jordan2024muon}
K.~Jordan, Y.~Jin, V.~Boza, J.~You, F.~Cesista, L.~Newhouse, and J.~Bernstein.
\newblock {Muon}: {A}n optimizer for hidden layers in neural networks.
\newblock {\em URL https://kellerjordan. github. io/posts/muon}, 6(3):4, 2024.

\bibitem{KingmaB14}
Diederik~P. Kingma and Jimmy Ba.
\newblock Adam: A method for stochastic optimization.
\newblock {\em CoRR}, abs/1412.6980, 2014.

\bibitem{Konnov2001}
I.V. Konnov.
\newblock {\em Combined relaxation methods for variational inequalities.}
\newblock Springer-Verlag, 2001.

\bibitem{Krasnoselskii1955}
M.~A. Krasnosel'ski\v{i}.
\newblock Two remarks on the method of successive approximations.
\newblock {\em Uspekhi Matematicheskikh Nauk}, 10(1):123--127, 1955.
\newblock In Russian.

\bibitem{lan2026optimal}
G.~Lan, Y.~Ouyang, and Z.~Zhang.
\newblock Optimal and parameter-free gradient minimization methods for convex
  and nonconvex optimization.
\newblock {\em Math. Program.}, pages 1--40, 2026.

\bibitem{li2025simple}
T.~Li and G.~Lan.
\newblock A simple uniformly optimal method without line search for convex
  optimization.
\newblock {\em Math. Program.}, pages 1--38, 2025.

\bibitem{lieder2021convergence}
F.~Lieder.
\newblock On the convergence rate of the halpern-iteration.
\newblock {\em Optim. Letters}, 15(2):405--418, 2021.

\bibitem{lv2026preconditioned}
F.~Lv and Q.-L. Dong.
\newblock Preconditioned {H}alpern iteration with adaptive anchoring parameters
  and an acceleration to {C}hambolle--{P}ock algorithm.
\newblock {\em Numerical Algorithms}, pages 1--30, 2026.

\bibitem{malitsky2019golden}
Y.~Malitsky.
\newblock Golden ratio algorithms for variational inequalities.
\newblock {\em Math. Program.}, 184(1--2):383--410, 2020.

\bibitem{malitsky2020adaptive}
Y.~Malitsky and K.~Mishchenko.
\newblock Adaptive gradient descent without descent.
\newblock In {\em International Conference on Machine Learning}, pages
  6702--6712. PMLR, 2020.

\bibitem{malitsky2016first}
Y.~Malitsky and T.~Pock.
\newblock A first-order primal-dual algorithm with linesearch.
\newblock {\em arXiv preprint arXiv:1608.08883}, 2016.

\bibitem{malitsky2020forward}
Y.~Malitsky and M.~K. Tam.
\newblock A forward-backward splitting method for monotone inclusions without
  cocoercivity.
\newblock {\em SIAM J. Optim.}, 30(2):1451--1472, 2020.

\bibitem{mann1953mean}
W.~R. Mann.
\newblock Mean value methods in iteration.
\newblock {\em Proceedings of the American Mathematical Society},
  4(3):506--510, 1953.

\bibitem{ogwo2025inertial}
G.~N. Ogwo, B.~Zinsou, H.~A. Abass, and O.~K. Oyewole.
\newblock Inertial {H}alpern-type {T}seng’s method for approximating a
  solution to monotone inclusion problems with fixed point constraint.
\newblock {\em Rendiconti del Circolo Matematico di Palermo Series 2},
  74(1):61, 2025.

\bibitem{orabona2016coin}
F.~Orabona and D.~P{\'a}l.
\newblock Coin betting and parameter-free online learning.
\newblock {\em Advances in Neural Information Processing Systems}, 29, 2016.

\bibitem{oyewole2022totally}
O.~K. Oyewole and S.~Reich.
\newblock A totally relaxed self-adaptive algorithm for solving a variational
  inequality and fixed point problems in banach spaces.
\newblock {\em Appl Set-Valued Anal Optim}, 4:349--366, 2022.

\bibitem{park2022exact}
J.~Park and E.~K. Ryu.
\newblock Exact optimal accelerated complexity for fixed-point iterations.
\newblock In {\em International Conference on Machine Learning}, pages
  17420--17457. PMLR, 2022.

\bibitem{sabach2017first}
S.~Sabach and S.~Shtern.
\newblock A first order method for solving convex bilevel optimization
  problems.
\newblock {\em SIAM J. Optim.}, 27(2):640--660, 2017.

\bibitem{shen2026parameter}
L.~Shen and F.~K{\i}l{\i}n{\c{c}}-Karzan.
\newblock Parameter-free non-ergodic extragradient algorithms for solving
  monotone variational inequalities.
\newblock {\em arXiv preprint arXiv:2604.07662}, 2026.

\bibitem{tan2022self}
B.~Tan and X.~Qin.
\newblock Self adaptive viscosity-type inertial extragradient algorithms for
  solving variational inequalities with applications.
\newblock {\em Mathematical Modelling and Analysis}, 27(1):41--58, 2022.

\bibitem{tran2022connection}
Q.~Tran-Dinh.
\newblock From {H}alpern's fixed-point iterations to {N}esterov's accelerated
  interpretations for root-finding problems.
\newblock {\em Comput. Optim. Appl.}, 87(1):181--218, 2024.

\end{thebibliography}

\end{document}

